%% Question bank chapter file - generated by bank_tools.py build (do not edit while the bank is still being built:
%% it is regenerated from the per-test data in Question Bank\data). Every question is one bankq block:
%%   \begin{bankq}[id=..., source={...}, paper=.., qno=.., type=..., subtopic={...}, answer={...}, answer_src=key|solved,
%%                 working={...}, difficulty=1-10, diff_parts={K S I C F}, diff_note={...},
%%                 group=<passage id or empty>, also_in={duplicate copies in other tests}]
%% difficulty = min(10, K+S+I+C+F-1): K concepts 1-3, S steps 1-4, I insight 0-2, C computation 0-2, F format 0-1
%%   ...question LaTeX, options, figures (\bfig / \bfigi -> img/)...
%%   \end{bankq}
%% Passages shared by several questions are bankpassage blocks; their questions carry group=<passage id>.
\documentclass[10pt]{article}
\usepackage[a4paper,margin=1.8cm]{geometry}
\usepackage[T1]{fontenc}
\usepackage{amsmath,amssymb,mathtools}
\usepackage{graphicx}
\usepackage[export]{adjustbox}
\usepackage{xcolor,needspace,keyval}
\usepackage{enumitem}
\usepackage{booktabs,array,tabularx}
\usepackage[hidelinks]{hyperref}
\setlength{\parindent}{0pt}
\setlength{\parskip}{3pt}
\makeatletter
\define@key{bq}{id}{\def\bq@id{#1}}\define@key{bq}{source}{\def\bq@src{#1}}\define@key{bq}{paper}{\def\bq@paper{#1}}
\define@key{bq}{qno}{\def\bq@qno{#1}}\define@key{bq}{type}{\def\bq@type{#1}}\define@key{bq}{subtopic}{\def\bq@sub{#1}}
\define@key{bq}{answer}{\def\bq@ans{#1}}\define@key{bq}{answer_src}{\def\bq@asrc{#1}}\define@key{bq}{working}{\def\bq@work{#1}}
\define@key{bq}{group}{\def\bq@group{#1}}\define@key{bq}{also_in}{\def\bq@also{#1}}
\define@key{bq}{difficulty}{\def\bq@diff{#1}}\define@key{bq}{diff_parts}{\def\bq@dparts{#1}}\define@key{bq}{diff_note}{\def\bq@dnote{#1}}
\newenvironment{bankq}[1][]{\def\bq@work{}\def\bq@also{}\def\bq@group{}\def\bq@diff{?}\def\bq@dparts{}\def\bq@dnote{}\setkeys{bq}{#1}%
  \par\needspace{4\baselineskip}\medskip\noindent{\small\textbf{\bq@id}\hfill\color{gray}[\bq@src, Paper \bq@paper, Q\bq@qno]
  \bq@type\ $\cdot$ \bq@sub\ $\cdot$ \textbf{D\bq@diff}}\par\nopagebreak\relax}
  {\par{\small\color{teal}Answer: \textbf{\bq@ans}\ifx\bq@work\@empty\else\ (solved: \bq@work)\fi
  \ifx\bq@also\@empty\else\par\color{gray}Also in: \bq@also\fi}\par}
\define@key{bp}{id}{\def\bp@id{#1}}\define@key{bp}{source}{\def\bp@src{#1}}\define@key{bp}{questions}{\def\bp@qs{#1}}
\newenvironment{bankpassage}[1][]{\setkeys{bp}{#1}\par\needspace{4\baselineskip}\bigskip\noindent
  {\small\textbf{Passage \bp@id}\hfill\color{gray}[\bp@src; questions \bp@qs]}\par\nopagebreak\relax\itshape}{\par}
\makeatother
\newcommand{\bfig}[2]{\par\smallskip{\centering\includegraphics[width=#1pt,max width=0.62\linewidth,max height=0.32\textheight]{img/#2}\par}\smallskip}
\newcommand{\bfigi}[2]{\includegraphics[width=#1pt,max width=\linewidth,max height=0.32\textheight,valign=c]{img/#2}}
\newcommand{\bmissing}[1]{\textit{[figure missing in the source conversion: \detokenize{#1}]}}
\newenvironment{qoptions}{\begin{itemize}[leftmargin=1.8em,itemsep=1pt,topsep=2pt,label={}]}{\end{itemize}}
\newcommand{\opt}[2]{\item[(#1)] #2}
\providecommand{\fwstart}{}\providecommand{\fwend}{}\providecommand{\qfullsplit}{}
\newcolumntype{L}[1]{>{\raggedright\arraybackslash}p{#1}}
\newcolumntype{C}[1]{>{\centering\arraybackslash}p{#1}}
\newcolumntype{R}[1]{>{\raggedleft\arraybackslash}p{#1}}
\providecommand{\sectionhdr}[1]{}\providecommand{\subsectionhdr}[1]{}


\begin{document}
{\Large\bfseries Salt Analysis --- Question Bank}\par\smallskip
{\small Chemistry chapter C33. Questions from all converted test papers, grouped by source series. Answers: printed key where it matches the paper, otherwise solved (answer\_src=solved, with working). Exact repeats across papers appear once, with also\_in listing the other copies.}\par\medskip\hrule

\section*{Score-I papers}
\begin{bankq}[id=SI12-P1-CHEM-S1-Q09, source={Score-I - Test 12 (09 Mar 2022)}, paper=1, qno=9, type=multi, subtopic={Cation tests}, answer={A,D}, answer_src=key, difficulty=2, diff_parts={K1 S1 I0 C0 F1}, diff_note={concepts: action of dilute HCl on sulphides | steps: HgS and CuS have too small Ksp to react}]
Which of the following sulphide(s) does/do not liberate H$_2$S on warning with dil. HCl.
\begin{qoptions}
\opt{A}{HgS}
\opt{B}{ZnS}
\opt{C}{FeS}
\opt{D}{CuS}
\end{qoptions}
\end{bankq}
\begin{bankq}[id=SI12-P2-CHEM-S2-Q06, source={Score-I - Test 12 (09 Mar 2022)}, paper=2, qno=6, type=numerical, subtopic={Cation tests}, answer={3}, answer_src=solved, working={Black sulphides dissolving in hot dil. HNO$_3$: CuS, Ag$_2$S, PbS (HgS needs aqua regia; Fe$^{3+}$ gives sulphur)}, difficulty=4, diff_parts={K1 S2 I1 C0 F1}, diff_note={concepts: qualitative analysis of group II cations | steps: black sulphides: CuS, PbS, Ag2S, HgS; HgS does not dissolve in hot dilute nitric acid; count 3 | I: HgS exception}]
Find out the number of cation(s) which form(s) black ppt (soluble in hot and dil HNO$_3$) on passing H$_2$S gas into their salt solution.\\
Mg$^{2+}_{(aq.)}$, Cu$^{2+}_{(aq.)}$, Ba$^{2+}_{(aq.)}$, Fe$^{3+}_{(aq.)}$, Ag$^+_{(aq.)}$, Al$^{3+}_{(aq.)}$, Hg$^{2+}_{(aq.)}$, Pb$^{2+}_{(aq.)}$, Mn$^{2+}_{(aq.)}$
\end{bankq}

\section*{Score-II papers}
\begin{bankq}[id=SII02-P1-CHEM-S3-Q02, source={Score-II - Test 02 (28 Mar 2021)}, paper=1, qno=2, type=numerical, subtopic={Anion tests}, answer={3.00}, answer_src=key, difficulty=8, diff_parts={K3 S1 I2 C2 F1}, diff_note={concepts: standard problem in this chapter, anion tests, apply the relevant results | steps: apply the relevant results; compute}]
A solution of white crystals gives a precipitate with AgNO$_3$ but no precipitation with a solution of Na$_2$CO$_3$. The action of concentrated H$_2$SO$_4$ on the crystals yields coloured fumes. Find the number of crystals given below which can give above observation:\\
NaI, KCl, NaNO$_2$, NaBr, CH$_3$COONa
\end{bankq}
\begin{bankq}[id=SII02-P2-CHEM-S2-Q11, source={Score-II - Test 02 (28 Mar 2021)}, paper=2, qno=11, type=single, subtopic={Anion tests}, answer={B}, answer_src=key, difficulty=4, diff_parts={K2 S1 I1 C1 F0}, diff_note={concepts: standard problem in this chapter, anion tests | steps: apply the relevant results; compute}]
Which combination will give X \& Y as colourless solution.

\begin{qoptions}
\opt{A}{(I) - (T)}
\opt{B}{(II) - (Q)}
\opt{C}{(III) - (P)}
\opt{D}{(IV) - (U)}
\end{qoptions}
\end{bankq}
\begin{bankq}[id=SII02-P2-CHEM-S2-Q12, source={Score-II - Test 02 (28 Mar 2021)}, paper=2, qno=12, type=single, subtopic={Anion tests}, answer={B}, answer_src=key, difficulty=4, diff_parts={K2 S1 I1 C1 F0}, diff_note={concepts: standard problem in this chapter, anion tests | steps: apply the relevant results; compute}]
Select INCORRECT combination

\begin{qoptions}
\opt{A}{(I) - (T)}
\opt{B}{(III) - (S)}
\opt{C}{(I) - (Q)}
\opt{D}{(IV) - (R)}
\end{qoptions}
\end{bankq}
\begin{bankq}[id=SII04-P1-CHEM-S1-Q07, source={Score-II - Test 04 (07 May 2021)}, paper=1, qno=7, type=single, subtopic={Anion tests}, answer={B}, answer_src=key, difficulty=4, diff_parts={K2 S1 I1 C1 F0}, diff_note={concepts: standard problem in this chapter, anion tests | steps: apply the relevant results; compute}]
Identify how many ``Ore names'', ``Type of Anion'', ``Metal(s) present in mineral'' are correctly related
\begin{center}
\renewcommand{\arraystretch}{1.15}
\begin{tabular}{@{}lll@{}}
(i) Magnetite & oxide & Fe\\
(ii) Cryolite & fluoride & Al\\
(iii) Copper Pyrite & Sulphide & Fe\\
(iv) Galena & sulphide & Hg\\
(v) Asbestos & Silicate & Al\\
(vi) Lunar Caustic & hydroxide & Ag\\
(vii) Dolomite & carbonate & Ca\\
(viii) Dolomite & hydroxide & Mg\\
\end{tabular}
\end{center}

\begin{qoptions}
\opt{A}{3}
\opt{B}{4}
\opt{C}{5}
\opt{D}{6}
\end{qoptions}
\end{bankq}
\begin{bankq}[id=SII04-P1-CHEM-S1-Q08, source={Score-II - Test 04 (07 May 2021)}, paper=1, qno=8, type=single, subtopic={Anion tests}, answer={C}, answer_src=key, difficulty=4, diff_parts={K2 S1 I1 C1 F0}, diff_note={concepts: standard problem in this chapter, anion tests | steps: apply the relevant results; compute}]
When a mixture of NaCl and K$_2$Cr$_2$O$_7$ is gently warmed with conc. H$_2$SO$_4$\\
(i) an orange-red vapour is evolved\\
(ii) the vapour formed when passed into NaOH solution gives yellow solution of Na$_2$CrO$_4$\\
(iii) Chlorine gas is not evolved at all\\
(iv) chromyl chloride gas is formed\\
(v) SnCl$_4$ used instead of NaCl also gives same observations

\begin{qoptions}
\opt{A}{All statements are correct}
\opt{B}{2 statements are correct}
\opt{C}{3 statements are correct}
\opt{D}{none of the statements are correct}
\end{qoptions}
\end{bankq}
\begin{bankq}[id=SII04-P2-CHEM-S4-Q05, source={Score-II - Test 04 (07 May 2021)}, paper=2, qno=5, type=numerical, subtopic={Anion tests}, answer={4}, answer_src=key, difficulty=8, diff_parts={K3 S1 I2 C2 F1}, diff_note={concepts: standard problem in this chapter, anion tests, apply the relevant results | steps: apply the relevant results; compute}]
Select the total number of substance which is/are given below existing as most stable allotrope of corresponding element and having their standard enthalpy of formation zero-\\
(i) Graphite \hspace{1em} (ii) Ozone \hspace{1em} (iii) Black phosphorus \hspace{1em} (iv) Rhombic sulphur\\
(v) White tin \hspace{1em} (vi) Cl$_2$(g) \hspace{1em} (vii) Br$_2$(g)
\end{bankq}
\begin{bankq}[id=SII06-P1-CHEM-S2-Q07, source={Score-II - Test 06 (05 Jul 2022)}, paper=1, qno=7, type=multi, subtopic={Anion tests}, answer={A,B,D}, answer_src=key, difficulty=8, diff_parts={K3 S1 I2 C2 F1}, diff_note={concepts: standard problem in this chapter, anion tests, apply the relevant results | steps: apply the relevant results; compute}]
Sometimes yellow turbidity appears while passing H$_2$S gas even in the absence of $2^{nd}$ group of basic radicals. Which of the following is/are not correct justification for the observation:
\begin{qoptions}
\opt{A}{Sulphur is present in mixture as impurity}
\opt{B}{IV group radicals are precipitated as sulphides.}
\opt{C}{H$_2$S is oxidised by some acid radicals}
\opt{D}{III$^{rd}$ group radicals are precipitated as yellow coloured hydroxides}
\end{qoptions}
\end{bankq}
\begin{bankq}[id=SII06-P1-CHEM-S3-Q01, source={Score-II - Test 06 (05 Jul 2022)}, paper=1, qno=1, type=numerical, subtopic={Anion tests}, answer={3}, answer_src=key, difficulty=8, diff_parts={K3 S1 I2 C2 F1}, diff_note={concepts: standard problem in this chapter, anion tests, apply the relevant results | steps: apply the relevant results; compute}]
How many of the following salts are white and evolves colourless gas or vapours with dil. H$_2$SO$_4$.\\
BaCrO$_4$, Hg$_2$CrO$_4$, ZnS, BaSO$_4$, BaS$_2$O$_3$, CH$_3$COOAg, AgNO$_2$.
\end{bankq}
\begin{bankq}[id=SII08-P1-CHEM-S2-Q10, source={Score-II - Test 08 (05 Feb 2023)}, paper=1, qno=10, type=single, subtopic={Anion tests}, answer={C}, answer_src=key, difficulty=4, diff_parts={K2 S1 I1 C1 F0}, diff_note={concepts: standard problem in this chapter, anion tests | steps: apply the relevant results; compute}]
Match the compounds in LIST-I with the observation in LIST-II, and choose the correct option.

\begin{center}
\begin{tabular}{|c|l|c|p{6cm}|}
\hline
\multicolumn{2}{|c|}{LIST-I} & \multicolumn{2}{c|}{LIST-II} \\
\hline
(I) & Aniline & (P) & Sodium fusion extract of the compound on boiling with FeSO$_4$, followed by acidification with conc. H$_2$SO$_4$, gives Prussian blue color. \\
\hline
(II) & o-Cresol & (Q) & Sodium fusion extract of the compound on treatment with sodium nitroprusside gives blood red color. \\
\hline
(III) & Cysteine & (R) & Addition of the compound to a saturated solution of NaHCO$_3$ results in effervescence. \\
\hline
(IV) & Coprolactam & (S) & The compound reacts with bromine water to give a white precipitate. \\
\hline
& & (T) & Treating the compound with neutral FeCl$_3$ solution produces violet color. \\
\hline
\end{tabular}
\end{center}

\begin{qoptions}
\opt{A}{I$\to$P; II$\to$R,S; III$\to$R; IV$\to$Q,S}
\opt{B}{I$\to$P,Q; II$\to$S; III$\to$Q,R; IV$\to$P}
\opt{C}{I$\to$P,S; II$\to$T; III$\to$Q,R; IV$\to$P}
\opt{D}{I$\to$Q,S; II$\to$P,T; III$\to$P; IV$\to$S}
\end{qoptions}
\end{bankq}
\begin{bankq}[id=SII08-P1-CHEM-S3-Q04, source={Score-II - Test 08 (05 Feb 2023)}, paper=1, qno=4, type=numerical, subtopic={Anion tests}, answer={0.225}, answer_src=key, difficulty=8, diff_parts={K3 S1 I2 C2 F1}, diff_note={concepts: standard problem in this chapter, anion tests, apply the relevant results | steps: apply the relevant results; compute}]
The treatment of an aqueous solution of 3.74 g of Cu(NO$_3$)$_2$ with excess KI results in a brown solution along with the formation of a precipitate. Passing H$_2$S through this brown solution gives another precipitate X. The precipitate 'X' on burning with excess of air produces a gas "Y". Volume of gas "Y" (in litre) at STP is\_\_\_\_\_\_\_.\\
\relax[Given: Atomic mass of H=1, N=14, O=16, S=32, K=39, Cu=63, I=127]
\end{bankq}
\begin{bankq}[id=SII08-P2-CHEM-S1-Q03, source={Score-II - Test 08 (05 Feb 2023)}, paper=2, qno=3, type=single, subtopic={Anion tests}, answer={C}, answer_src=key, difficulty=4, diff_parts={K2 S1 I1 C1 F0}, diff_note={concepts: standard problem in this chapter, anion tests | steps: apply the relevant results; compute}]
The reaction Pb(NO$_3$)$_2$ and NaCl in water produces a precipitate that dissolves upon the addition of HCl of appropriate concentration. The dissolution of the precipitate is due to the formation of
\begin{qoptions}
\opt{A}{PbCl$_2$}
\opt{B}{PbCl$_4$}
\opt{C}{$[\mathrm{PbCl_4}]^{2-}$}
\opt{D}{$[\mathrm{PbCl_6}]^{2-}$}
\end{qoptions}
\end{bankq}
\begin{bankq}[id=SII09-P2-CHEM-S1-Q02, source={Score-II - Test 09 (12 Feb 2023)}, paper=2, qno=2, type=single, subtopic={Anion tests}, answer={B}, answer_src=key, difficulty=4, diff_parts={K2 S1 I1 C1 F0}, diff_note={concepts: standard problem in this chapter, anion tests | steps: apply the relevant results; compute}]
'S' is mixture of two salts. Consider the following facts about this mixture:\\
a. Aqueous solution of S $\xrightarrow{H_2SO_4}$ White precipitate (W) + Coloured Solution (C)\\
b. W + $CH_3COONH_4$ (conc) $\to$ Soluble complex\\
Appropriate option(s) for the mixture of the salt(s) is/are:
\begin{qoptions}
\opt{A}{$AgNO_3$ and $Zn(NO_3)_2$}
\opt{B}{$Pb(NO_3)_2$ and $Cu(NO_3)_2$}
\opt{C}{$Zn(NO_3)_2$ and $Cu(NO_3)_2$}
\opt{D}{$Pb(NO_3)_2$ and $Zn(NO_3)_2$}
\end{qoptions}
\end{bankq}
\begin{bankpassage}[id=SII14-P1-CHEM-S3-G3-4, source={Score-II - Test 14 (06 Feb 2024)}, questions={3-4}]
An aqueous blue coloured solution of a transition metal sulphate reacts with $H_2S$ in acidic medium to give a black precipitate A, which is insoluble in warm aqueous solution of KOH. The blue solution on treatment with KI in weakly acidic medium, turns yellow / brown solution due to the formations of compound C and produces a white precipitate B.\\
(given Atomic mass of Cu = 63.5, S = 32, O = 16, N = 14, I = 127, K = 39, H = 1)
\end{bankpassage}
\begin{bankq}[id=SII14-P1-CHEM-S3-Q03, source={Score-II - Test 14 (06 Feb 2024)}, paper=1, qno=3, type=numerical, subtopic={Anion tests}, answer={4.00}, answer_src=key, difficulty=8, diff_parts={K3 S1 I2 C2 F1}, diff_note={concepts: standard problem in this chapter, anion tests, apply the relevant results | steps: apply the relevant results; compute}, group=SII14-P1-CHEM-S3-G3-4]
Number of atoms in per formula unit of C.
\end{bankq}
\begin{bankq}[id=SII14-P1-CHEM-S3-Q04, source={Score-II - Test 14 (06 Feb 2024)}, paper=1, qno=4, type=numerical, subtopic={Anion tests}, answer={168.00}, answer_src=key, difficulty=8, diff_parts={K3 S1 I2 C2 F1}, diff_note={concepts: standard problem in this chapter, anion tests, apply the relevant results | steps: apply the relevant results; compute}, group=SII14-P1-CHEM-S3-G3-4]
"A" when dissolve in dilute nitric acid white cloudy appearance produce in test tube, then how many grams of dilute $HNO_3$ is required to dissolve 1 mole of "A" ?
\end{bankq}
\begin{bankq}[id=SII16-P1-CHEM-S4-Q09, source={Score-II - Test 16 (09 Feb 2024)}, paper=1, qno=9, type=numerical, subtopic={Anion tests}, answer={12}, answer_src=key, difficulty=8, diff_parts={K3 S1 I2 C2 F1}, diff_note={concepts: standard problem in this chapter, anion tests, apply the relevant results | steps: apply the relevant results; compute}]
A black colored compound (A) on reaction with dilute $H_2SO_4$ gives a gas (B) which on passing in a solution of an acid (C) gives a white turbidity (D). Gas (B) when passed in an acidified solution of a compound (E) gives a precipitate (F) soluble in dilute $HNO_3$. After boiling this solution when an excess of $NH_4OH$ is added, a blue colored compound (G) is formed. To this solution on addition of acetic acid and aqueous $K_4[Fe(CN)_6]$ a chocolate precipitate (H) is obtained. On addition of an aqueous solution of $BaCl_2$ to an aqueous solution of (E), a white precipitate insoluble in $HNO_3$ is obtained.\\
Calculate the sum of oxidation state of every different type of metal present in the compounds from A to H.
\end{bankq}
\begin{bankq}[id=SII18-P1-CHEM-S3-Q03, source={Score-II - Test 18 (13 Feb 2024)}, paper=1, qno=3, type=numerical, subtopic={Anion tests}, answer={Bonus}, answer_src=solved, working={Printed key page is incomplete/misaligned for this question and the question was not solved; banked as Bonus.}, difficulty=8, diff_parts={K3 S1 I2 C2 F1}, diff_note={concepts: standard problem in this chapter, anion tests, apply the relevant results | steps: apply the relevant results; compute}]
Number of following salts which is/are water insoluble but neither black nor white nor blue in color, is X then value of 2X.

Cu(CN)$_2$, Cu(OH)$_2$, Cu$_2$O, Cu$_3$[Fe(CN)$_6$]$_2$, Co(OH)NO$_3$, HgO.HgNH$_2$I, HgNH$_2$Cl, CuSO$_4$, Ag$_3$[Fe(CN)$_6$]
\end{bankq}
\begin{bankq}[id=SII21-P1-CHEM-S2-Q10, source={Score-II - Test 21 (14 Mar 2024)}, paper=1, qno=10, type=single, subtopic={Anion tests}, answer={D}, answer_src=key, difficulty=4, diff_parts={K2 S1 I1 C1 F0}, diff_note={concepts: standard problem in this chapter, anion tests | steps: apply the relevant results; compute}]
Match the compounds in LIST-I with the observation in LIST-II, and choose the correct option.

\begin{center}
\begin{tabular}{|c|c|c|p{6.5cm}|}
\hline
\multicolumn{2}{|c|}{LIST-I} & \multicolumn{2}{c|}{LIST-II} \\
\hline
(I) & Aniline & (P) & Sodium fusion extract of the compound on boiling with FeSO$_4$, followed by acidification with conc. H$_2$SO$_4$, gives Prussian blue color. \\
\hline
(II) & o-Cresol & (Q) & Sodium fusion extract of the compound on treatment with sodium nitroprusside gives blood red color. \\
\hline
(III) & Cysteine & (R) & Addition of the compound to a saturated solution of NaHCO$_3$ results in effervescence. \\
\hline
(IV) & Coprolactam & (S) & The compound reacts with bromine water to give a white precipitate. \\
\hline
& & (T) & Treating the compound with neutral FeCl$_3$ solution produces violet color. \\
\hline
\end{tabular}
\end{center}

\begin{qoptions}
\opt{A}{I$\to$P,Q; II$\to$S; III$\to$Q,R; IV$\to$P}
\opt{B}{I$\to$P; II$\to$R,S; III$\to$R; IV$\to$Q,S}
\opt{C}{I$\to$Q,S; II$\to$P,T; III$\to$P; IV$\to$S}
\opt{D}{I$\to$P,S; II$\to$T; III$\to$Q,R; IV$\to$P}
\end{qoptions}
\end{bankq}
\begin{bankq}[id=SII21-P1-CHEM-S3-Q04, source={Score-II - Test 21 (14 Mar 2024)}, paper=1, qno=4, type=numerical, subtopic={Anion tests}, answer={Bonus}, answer_src=solved, working={Printed key omits or corrupts this numeric answer and the question was not solved; banked as Bonus.}, difficulty=8, diff_parts={K3 S1 I2 C2 F1}, diff_note={concepts: standard problem in this chapter, anion tests, apply the relevant results | steps: apply the relevant results; compute}]
The treatment of an aqueous solution of 7.48 g of Cu(NO$_3$)$_2$ with excess KI results in a brown solution along with the formation of a precipitate. Passing H$_2$S through this brown solution gives another precipitate X. The amount of X (in g) is\underline{\hspace{2em}}.
[Given : Atomic mass of H = 1, N = 14, O = 16, S = 32, K = 39, Cu = 63, I = 127]
\end{bankq}
\begin{bankq}[id=SII21-P2-CHEM-S1-Q03, source={Score-II - Test 21 (14 Mar 2024)}, paper=2, qno=3, type=single, subtopic={Anion tests}, answer={D}, answer_src=key, difficulty=4, diff_parts={K2 S1 I1 C1 F0}, diff_note={concepts: standard problem in this chapter, anion tests | steps: apply the relevant results; compute}]
The reaction Pb(NO$_3$)$_2$ and NaCl in water produces a precipitate that dissolves upon the addition of HCl of appropriate concentration. The dissolution of the precipitate is due to the formation of
\begin{qoptions}
\opt{A}{PbCl$_2$}
\opt{B}{PbCl$_4$}
\opt{C}{[PbCl$_6$]$^{2-}$}
\opt{D}{[PbCl$_4$]$^{2-}$}
\end{qoptions}
\end{bankq}
\begin{bankq}[id=SII26-P1-CHEM-S2-Q08, source={Score-II - Test 26 (20 Mar 2024)}, paper=1, qno=8, type=multi, subtopic={Anion tests}, answer={A}, answer_src=key, difficulty=5, diff_parts={K2 S1 I1 C1 F1}, diff_note={concepts: standard problem in this chapter, anion tests | steps: apply the relevant results; compute}]
\bfig{130.3}{SII26_p1_chem_s2_q08_1.png}
`B' + Ca(OH)$_2$(aq) $\longrightarrow$ solution turns milky\\
`B' gives green solution with K$_2$Cr$_2$O$_7$(aq) in acidic medium.\\
Find correct statement(s)
\begin{qoptions}
\opt{A}{`C' has structure similar to H$_2$O$_2$}
\opt{B}{`B' is SO$_3$}
\opt{C}{`E' forms a complex with CuSO$_4$(aq) solution}
\opt{D}{`D' is dibasic acid}
\end{qoptions}
\end{bankq}
\begin{bankq}[id=SII32-P1-CHEM-S2-Q05, source={Score-II - Test 32 (29 Mar 2024)}, paper=1, qno=5, type=multi, subtopic={Anion tests}, answer={B,C}, answer_src=key, difficulty=8, diff_parts={K3 S1 I2 C2 F1}, diff_note={concepts: standard problem in this chapter, anion tests, apply the relevant results | steps: apply the relevant results; compute}]
A salt (white) is treated with water to form a blue color solution, which on treatment with potassium hexacyanoferrate(II) forms a chocolate brown precipitate 'X':

\begin{center}
\bfigi{233.8}{SII32_chem_q5_reaction.png}
\end{center}

if ppt.\ of 'x' is soluble in ammonia solution, and forms dark-blue solution of 'y' which contains two complex ion. Which of the following option is/are correct for 'y' solution?

\begin{qoptions}
\opt{A}{'y' is the solution of [Cu(CN)$_4$]$^{-2}$ and [Fe(CN)$_6$]$^{-3}$ ions.}
\opt{B}{One of the complex can show structural isomerism.}
\opt{C}{One of the complex ion in solution of 'y' produces blue ppt.\ with FeCl$_3$ solution.}
\opt{D}{One of the complex has tetrahedral geometry and other complex has octahedral geometry.}
\end{qoptions}
\end{bankq}
\begin{bankq}[id=SII32-P1-CHEM-S3-Q01, source={Score-II - Test 32 (29 Mar 2024)}, paper=1, qno=1, type=numerical, subtopic={Anion tests}, answer={3.24}, answer_src=key, difficulty=8, diff_parts={K3 S1 I2 C2 F1}, diff_note={concepts: standard problem in this chapter, anion tests, apply the relevant results | steps: apply the relevant results; compute}]
The value of \textbf{x} is\_\_\_\_\_.
\end{bankq}
\begin{bankq}[id=SII32-P1-CHEM-S3-Q02, source={Score-II - Test 32 (29 Mar 2024)}, paper=1, qno=2, type=numerical, subtopic={Anion tests}, answer={3.20}, answer_src=key, difficulty=8, diff_parts={K3 S1 I2 C2 F1}, diff_note={concepts: standard problem in this chapter, anion tests, apply the relevant results | steps: apply the relevant results; compute}]
The value of \textbf{y} is\_\_\_\_\_.

\textbf{Paragraph for Questions 3 and 4}

A black salt (A) is insoluble in dilute acid in cold but dissolve in boiling solution of concentrated hydrochloric acid, evolving a pungent smelling gas B. B turned a filter paper moistened with lead acetate, black. The above solution was cooled and separated into two parts. One part was treated with NaOH solution giving a white precipitate (C), insoluble in excess of reagent. C, when added to a concentrated H$_2$O$_2$ solution, a clear, yellowish-brown solution of D was formed. Other part of the original solution was treated with KI solution, giving a black precipitate (E), which dissolved in excess of reagent forming a clear orange solution.
\end{bankq}
\begin{bankq}[id=SII32-P1-CHEM-S3-Q04, source={Score-II - Test 32 (29 Mar 2024)}, paper=1, qno=4, type=numerical, subtopic={Anion tests}, answer={5.00}, answer_src=key, difficulty=8, diff_parts={K3 S1 I2 C2 F1}, diff_note={concepts: standard problem in this chapter, anion tests, apply the relevant results | steps: apply the relevant results; compute}]
Given the reaction scheme:

\begin{center}
\bfigi{233.8}{SII32_chem_p34_q4_scheme.png}
\end{center}

Number of process which produce colored ppt./turbidity/solution (except white/colorless).
\end{bankq}
\begin{bankq}[id=SII44-P1-CHEM-S4-Q04, source={Score-II - Test 44 (23 Mar 2025)}, paper=1, qno=4, type=numerical, subtopic={Anion tests}, answer={72}, answer_src=key, difficulty=8, diff_parts={K3 S1 I2 C2 F1}, diff_note={concepts: standard problem in this chapter, anion tests, apply the relevant results | steps: apply the relevant results; compute}]
MgC$\ell_{2(aq)}$ $\xrightarrow[\text{NH}_4\text{C}\ell/\text{NH}_4\text{OH}]{\text{NaH}_2\text{PO}_4}$ (X)$\downarrow$ $\xrightarrow\Delta$ (Y)$\downarrow$ + other products \quad [(X): white ppt, (Y): white ppt]

Determine P$^{(\text{solubility of X})}$ at pH=10 and 25$^\circ$C. (Answer to the nearest integer)

Given: 'P' means $(-\log_{10})$

At equilibrium, $[\text{NH}_4^+]=0.01$M

K$_{sp}$(X) $=10^{-12}$ at 25$^\circ$C

For H$_3$PO$_4$: P$^{ka_1}=3$, P$^{ka_2}=7$, P$^{ka_3}=12$
\end{bankq}
\begin{bankq}[id=SII46-P1-CHEM-S2-Q05, source={Score-II - Test 46 (26 Mar 2025)}, paper=1, qno=5, type=numerical, subtopic={Anion tests}, answer={Bonus}, answer_src=solved, working={Printed key omits the numerical-section answers for this subject and the question was not solved; banked as Bonus.}, difficulty=8, diff_parts={K3 S1 I2 C2 F1}, diff_note={concepts: standard problem in this chapter, anion tests, apply the relevant results | steps: apply the relevant results; compute}]
How many of the following ions will form white precipitate with K$_4$[Fe(CN)$_6$] (before aerial oxidation)?
\[\mathrm{Co^{2+},\ Ca^{2+},\ Fe^{2+},\ Fe^{3+},\ Mn^{2+},\ Ag^+,\ Zn^{2+}}\]
\end{bankq}

\section*{Score Advanced papers}
\begin{bankq}[id=SA01-P1-CHEM-S1-Q04, source={Score Advanced - Test 01 (22 Sep 2020)}, paper=1, qno=4, type=single, subtopic={Anion tests}, answer={B}, answer_src=key, difficulty=2, diff_parts={K1 S1 I1 C0 F0}, diff_note={concepts: ferrocyanide with concentrated acid | steps: products of decomposition; select}]
Select the correct statement when potassium ferrocyanide is treated with concentrated sulphuric acid ?
\begin{qoptions}
\opt{A}{Fe$^{2+}$ of $\left[\mathrm{Fe(CN)_6}\right]^{4-}$ is oxidised to $\mathrm{FeO_2^-}$}
\opt{B}{CO gas is evolved}
\opt{C}{HCOOH is one of final product}
\opt{D}{HCN gas is one of the final products}
\end{qoptions}
\end{bankq}
\begin{bankq}[id=SA01-P2-CHEM-S1-Q01, source={Score Advanced - Test 01 (22 Sep 2020)}, paper=2, qno=1, type=multi, subtopic={Anion tests}, answer={A,B,C,D}, answer_src=key, difficulty=2, diff_parts={K1 S1 I0 C0 F1}, diff_note={concepts: gases from salts with dilute acid | steps: reactions of each salt}]
Which of the following produce(s) gaseous product on treatment with dilute HCl ?
\begin{qoptions}
\opt{A}{Na$_2$S$_2$O$_3$}
\opt{B}{NaHCO$_3$}
\opt{C}{Na$_2$SO$_3$}
\opt{D}{Na$_2$S}
\end{qoptions}
\end{bankq}
\begin{bankq}[id=SA06-P1-CHEM-S3-Q04, source={Score Advanced - Test 06 (04 Sep 2021)}, paper=1, qno=4, type=numerical, subtopic={Cation tests}, answer={3}, answer_src=key, difficulty=5, diff_parts={K2 S1 I1 C1 F1}, diff_note={concepts: precipitation with chrome alum, cation tests | steps: Cr3+ precipitates as hydroxide and sulphide, with carbonate}]
How many of the following compound(s) give precipitate with solution of chrome alum.\par
(i) BaCl$_2$ \quad (ii) Na$_2$CO$_3$ \quad (iii) (NH$_4$)$_2$S \quad (iv) Excess NaOH
\end{bankq}
\begin{bankq}[id=SA06-P2-CHEM-S2-Q08, source={Score Advanced - Test 06 (04 Sep 2021)}, paper=2, qno=8, type=single, subtopic={Cation tests}, answer={A}, answer_src=key, difficulty=7, diff_parts={K3 S1 I2 C2 F0}, diff_note={concepts: matching ions with their tests, cation tests, precipitation behaviour of Fe, Cu, Hg and Pb ions | steps: precipitation behaviour of Fe, Cu, Hg and Pb ions}]
{\footnotesize\setlength{\tabcolsep}{2pt}\noindent\begin{tabular}{|L{.6cm}|L{1cm}|L{.6cm}|L{3.8cm}|}
\hline
\multicolumn{2}{|c|}{List-I} & \multicolumn{2}{c|}{List-II} \\ \hline
(P) & Fe$^{n+}$ & (1) & (For n = 0), forms white ppt with dil. HCl. \\ \hline
(Q) & Cu$^{n+}$ & (2) & (For n = 2), forms coloured ppt with excess of NaOH solution. \\ \hline
(R) & Hg$^{n+}$ & (3) & (For n = 0), N$_2$O gas will be released with dil.HNO$_3$. \\ \hline
(S) & Pb$^{n+}$ & (4) & (For n = 2), Forms soluble complex with excess of NH$_3$ solution. \\ \hline
 & & (5) & (For n = 2) forms black ppt on passing H$_2$S gas in its aqueous solution. \\ \hline
\end{tabular}}\par
\begin{qoptions}
\opt{A}{P $\rightarrow$ 2,3;Q $\rightarrow$ 2,4,5;R $\rightarrow$ 2,5;S $\rightarrow$ 1,5}
\opt{B}{P $\rightarrow$ 1,2,3;Q $\rightarrow$ 3,5;R $\rightarrow$ 1,2,5;S $\rightarrow$ 2,4,5}
\opt{C}{P $\rightarrow$ 1,4;Q $\rightarrow$ 2,4;R $\rightarrow$ 3,5;S $\rightarrow$ 1,3}
\opt{D}{P $\rightarrow$ 2,4,5;Q $\rightarrow$ 2,3,5;R $\rightarrow$ 1,5;S $\rightarrow$ 1,4,5}
\end{qoptions}
\end{bankq}
\begin{bankq}[id=SA07-P1-CHEM-S3-Q01, source={Score Advanced - Test 07 (16 Apr 2023)}, paper=1, qno=1, type=numerical, subtopic={Cation tests}, answer={5}, answer_src=key, difficulty=5, diff_parts={K2 S1 I1 C1 F1}, diff_note={concepts: sulphides precipitated by H2S in acid, cation tests | steps: groups II sulphides}]
Among the given molecules, the number of metal sulphides precipitated by passing H$_2$S in acidic medium are ..........\par
ZnS, PbS, CuS, NiS, MnS, CdS, CoS, HgS, SnS
\end{bankq}
\begin{bankq}[id=SA07-P2-CHEM-S2-Q08, source={Score Advanced - Test 07 (16 Apr 2023)}, paper=2, qno=8, type=single, subtopic={Cation tests}, answer={B}, answer_src=key, difficulty=7, diff_parts={K3 S1 I2 C2 F0}, diff_note={concepts: reactions of reagents with solutions, cation tests, Na2O2, lead acetate, carbonate and iodate with the listed solutions | steps: Na2O2, lead acetate, carbonate and iodate with the listed solutions}]
\par
{\scriptsize\setlength{\tabcolsep}{1.5pt}\noindent\begin{tabular}{|L{.6cm}|L{3cm}|L{.6cm}|L{3.2cm}|}
\hline
 & List-I\par (Reagent) & & List-II\par (Observable change with given compounds) \\ \hline
(P) & Na$_2$O$_2$ solution & (1) & FeCl$_3$ solution \\ \hline
(Q) & Pb(CH$_3$COO)$_2$ solution & (2) & (NH$_4$)$_2$S solution \\ \hline
(R) & Na$_2$CO$_3$ solution & (3) & CuSO$_4$ solution \\ \hline
(S) & IO$_3^-$(aq) + H$^+$(aq) & (4) & KI solution \\ \hline
 & & (5) & Cr$_2$(SO$_4$)$_3$ solution \\ \hline
\end{tabular}}\par
\begin{qoptions}
\opt{A}{P $\rightarrow$ 2,3;Q $\rightarrow$ 1,2,3;R $\rightarrow$ 3,5;S $\rightarrow$ 2,3}
\opt{B}{P $\rightarrow$ 1,2,3,4,5;Q $\rightarrow$ 1,2,3,4,5;R $\rightarrow$ 1,2,3,5;S $\rightarrow$ 2,4}
\opt{C}{P $\rightarrow$ 3,4,5;Q $\rightarrow$ 2,3;R $\rightarrow$ 1,2;S $\rightarrow$ 1}
\opt{D}{P $\rightarrow$ 3,4;Q $\rightarrow$ 2,3;R $\rightarrow$ 1,2;S $\rightarrow$ 4,5}
\end{qoptions}
\end{bankq}
\begin{bankq}[id=SA09-P1-CHEM-S1-Q01, source={Score Advanced - Test 09 (29 Apr 2023)}, paper=1, qno=1, type=single, subtopic={Cation tests}, answer={C}, answer_src=key, difficulty=4, diff_parts={K2 S1 I1 C1 F0}, diff_note={concepts: identifying a coloured ammonium salt, cation tests | steps: NaOH liberates NH3; the chromate solution gives yellow BaCrO4}]
The coloured salt P was warmed with NaOH solution and colourless gas Q was evolved and a coloured solution R was produced. The colour of salt P was different from that of solution R. On adding BaCl$_2$ to R, a coloured PPT was obtained. Then the correct formula for salt P is
\begin{qoptions}
\opt{A}{Na$_2$Cr$_2$O$_7$}
\opt{B}{Mg(NH$_4$)PO$_4$}
\opt{C}{(NH$_4$)$_2$Cr$_2$O$_7$}
\opt{D}{PH$_4$I}
\end{qoptions}
\end{bankq}
\begin{bankpassage}[id=SA09-P2-CHEM-S2-G11-12, source={Score Advanced - Test 09 (29 Apr 2023)}, questions={11-12}]
Answer the following by appropriately matching the lists based on the information given in the paragraph.\par
{\footnotesize\setlength{\tabcolsep}{1.5pt}\noindent\begin{tabular}{|l|L{2.2cm}|l|L{3.4cm}|}
\hline
 & \multicolumn{1}{c|}{List - I} & & \multicolumn{1}{c|}{List - II}\\ \hline
(I) & K$_4$[Fe(CN)$_6$] & (P) & Produces deep blue solution with CuSO$_4$ solution\\ \hline
(II) & (NH$_4$)$_2$S & (Q) & Produces deep blue PPT with FeSO$_4$ solution\\ \hline
(III) & NH$_4$OH (excess) & (R) & Produces deep blue PPT with FeCl$_3$ solution\\ \hline
(IV) & NaOH (excess) & (S) & Produces colourless solution or white PPT with ZnSO$_4$ solution\\ \hline
 & & (T) & Produces chocolate brown PPT with CuSO$_4$ solution\\ \hline
 & & (U) & Produces reddish brown precipitate with FeCl$_3$\\ \hline
\end{tabular}}\par
\end{bankpassage}
\begin{bankq}[id=SA09-P2-CHEM-S2-Q11, source={Score Advanced - Test 09 (29 Apr 2023)}, paper=2, qno=11, type=single, subtopic={Cation tests}, answer={C}, answer_src=key, difficulty=7, diff_parts={K3 S1 I2 C2 F0}, diff_note={concepts: reactions of hexacyanoferrate(II) with cations, cation tests, Fe(III) gives Prussian blue | steps: Fe(III) gives Prussian blue; Cu(II) gives chocolate brown}, group=SA09-P2-CHEM-S2-G11-12]
Which of the following is incorrect match?
\begin{qoptions}
\opt{A}{I $\longrightarrow$ R}
\opt{B}{I $\longrightarrow$ S}
\opt{C}{I $\longrightarrow$ Q}
\opt{D}{I $\longrightarrow$ T}
\end{qoptions}
\end{bankq}
\begin{bankq}[id=SA09-P2-CHEM-S2-Q12, source={Score Advanced - Test 09 (29 Apr 2023)}, paper=2, qno=12, type=single, subtopic={Cation tests}, answer={D}, answer_src=key, difficulty=7, diff_parts={K3 S1 I2 C2 F0}, diff_note={concepts: reactions of ammonium sulphide and bases with cations, cation tests, zinc hydroxide dissolves in excess | steps: zinc hydroxide dissolves in excess; iron(III) gives reddish brown ppt}, group=SA09-P2-CHEM-S2-G11-12]
Which of the following is correct match?
\begin{qoptions}
\opt{A}{IV $\longrightarrow$ S,U}
\opt{B}{III $\longrightarrow$ P,S,U}
\opt{C}{II $\longrightarrow$ S}
\opt{D}{All of the above}
\end{qoptions}
\end{bankq}
\begin{bankq}[id=SA10-P1-CHEM-S2-Q10, source={Score Advanced - Test 10 (06 May 2023)}, paper=1, qno=10, type=single, subtopic={Cation tests}, answer={C}, answer_src=key, difficulty=5, diff_parts={K2 S2 I1 C1 F0}, diff_note={concepts: dissolution of sulphide and hydroxide precipitates, cation tests | steps: HgS dissolves in sodium sulphide; Ni(OH)2 in ammonia; MnS in acetic acid; Cr(OH)3 in alkali}]
List-I contains precipitate and List-II contains reagent in which it dissolves.\par
{\footnotesize\setlength{\tabcolsep}{1.5pt}\noindent\begin{tabular}{|c|l|c|l|}
\hline
 & List-I & & List-II\\ \hline
(P) & HgS & (1) & dil. HNO$_3$.\\ \hline
(Q) & Ni(OH)$_2$ & (2) & NaOH solution.\\ \hline
(R) & MnS & (3) & NH$_3$ solution.\\ \hline
(S) & Cr(OH)$_3$ & (4) & Na$_2$S (2M).\\ \hline
 & & (5) & Acetic acid solution.\\ \hline
\end{tabular}}\par
\begin{qoptions}
\opt{A}{P $\longrightarrow$ 1;Q $\longrightarrow$ 3;R $\longrightarrow$ 5;S $\longrightarrow$ 2}
\opt{B}{P $\longrightarrow$ 2;Q $\longrightarrow$ 1;R $\longrightarrow$ 3;S $\longrightarrow$ 4}
\opt{C}{P $\longrightarrow$ 4;Q $\longrightarrow$ 3;R $\longrightarrow$ 5;S $\longrightarrow$ 2}
\opt{D}{P $\longrightarrow$ 4;Q $\longrightarrow$ 2;R $\longrightarrow$ 1;S $\longrightarrow$ 5}
\end{qoptions}
\end{bankq}
\begin{bankq}[id=SA12-P1-CHEM-S3-Q03, source={Score Advanced - Test 12 (11 May 2023)}, paper=1, qno=3, type=numerical, subtopic={Cation tests}, answer={4}, answer_src=key, difficulty=5, diff_parts={K2 S1 I1 C1 F1}, diff_note={concepts: coordination number of mercury in the blue precipitate, cation tests | steps: the complex is Co[Hg(SCN)4]}, also_in={Score Advanced - Test 29 (01 May 2024) P2 Q3}]
A solution of Hg$^{2+}$ ion on treatment with a solution of cobalt (II) thiocyanate gives rise to a deep blue crystalline precipitate. The coordination number of mercury in the deep blue coloured compound is
\end{bankq}
\begin{bankq}[id=SA13-P1-CHEM-S2-Q04, source={Score Advanced - Test 13 (13 May 2023)}, paper=1, qno=4, type=single, subtopic={Anion tests}, answer={7}, answer_src=key, difficulty=7, diff_parts={K3 S1 I2 C2 F0}, diff_note={concepts: unknown compound giving gases with reagents, anion tests, identify the gases | steps: identify the gases; titration of the basic gas}]
\par
\bfig{130.5}{SA13_p1_chem_s2_q04_1.png}
Unknown mass of compound (P) reacted with different reagents to form gases (A), (B) and (C). Most basic gas released in above process was compltely titrated by 100 ml of 0.2M HCl.\par
If W$_\text{A}$ = mass of gas A (gas), in gm\par
W$_\text{B}$ = mass of gas B (gas), in gm\par
W$_\text{C}$ = mass of gas C (gas), in gm\par
Calculate : (W$_\text{A}$ + W$_\text{B}$ + W$_\text{C}$) $\times$ 10\par
\end{bankq}
\begin{bankpassage}[id=SA13-P2-CHEM-S1-G3-4, source={Score Advanced - Test 13 (13 May 2023)}, questions={3-4}]
Sodium salt of Pseudo halide X, Y and Z form colourless solution in water. Solution of X decolorizes I$_3^-$ solution in the presence of Y and gives brisk effervescence. Solution of Y gives an intense red colour compound ``R'' on reaction with Fe$^{+3}$ solution. Solution of Z gives intense blue colour ``B'' on reaction with a solution containing Fe$^{+2}$ \& Fe$^{+3}$ ions.\par
\end{bankpassage}
\begin{bankq}[id=SA13-P2-CHEM-S1-Q03, source={Score Advanced - Test 13 (13 May 2023)}, paper=2, qno=3, type=multi, subtopic={Anion tests}, answer={A,C}, answer_src=key, difficulty=8, diff_parts={K3 S1 I2 C2 F1}, diff_note={concepts: identifying pseudohalide ions from reactions, anion tests, thiocyanate, cyanate and azide behaviour | steps: thiocyanate, cyanate and azide behaviour}, group=SA13-P2-CHEM-S1-G3-4]
Select correct for X, Y and Z?
\begin{qoptions}
\opt{A}{anions in X, Y \& Z are nitrogen containing}
\opt{B}{anions in X, Y \& Z all have two $\sigma$ bonds}
\opt{C}{reaction of ``X'' and I$_3^-$ is a redox change}
\opt{D}{Colourless solution is obtained by reaction of I$_3^-$ \& X, produces violet blue solution with starch}
\end{qoptions}
\end{bankq}
\begin{bankq}[id=SA13-P2-CHEM-S1-Q04, source={Score Advanced - Test 13 (13 May 2023)}, paper=2, qno=4, type=multi, subtopic={Anion tests}, answer={A,B,C,D}, answer_src=key, difficulty=8, diff_parts={K3 S1 I2 C2 F1}, diff_note={concepts: further properties of the compounds R and B, anion tests, all follow from the identification | steps: all follow from the identification}, group=SA13-P2-CHEM-S1-G3-4]
Select correct for compound ``R'' and ``B''.
\begin{qoptions}
\opt{A}{``B'' is NaFe$^\text{III}$[Fe$^\text{II}$(CN)$_6$]}
\opt{B}{``R'' is [Fe(SCN)]$^{+2}$}
\opt{C}{``R'' is decolorized with excess KF$_\text{aq.}$}
\opt{D}{``R'' have $\sqrt{35}$ B.M. spin magnetic moment}
\end{qoptions}
\end{bankq}
\begin{bankq}[id=SA13-P2-CHEM-S3-Q13, source={Score Advanced - Test 13 (13 May 2023)}, paper=2, qno=13, type=single, subtopic={Cation tests}, answer={C}, answer_src=key, difficulty=7, diff_parts={K3 S1 I2 C2 F0}, diff_note={concepts: observations at each step of a flow chart, cation tests, match each step to its observation | steps: match each step to its observation}]
\textit{(Paragraph for Question 13)}\par
\bfig{332.3}{SA13_p2_chem_s3_q12_1.png}

{\footnotesize\setlength{\tabcolsep}{1.5pt}\noindent\begin{tabular}{|L{2.2cm}|c|L{4.3cm}|}
\hline
\textbf{List-I} (Steps in above flow chart) & & \textbf{List-II} (Observation by given step) \\
\hline
(a) \textbf{STEP-1} & (P) & Cu$^{+2}$ is separated from Hg$^{+2}$ \\
\hline
(b) \textbf{STEP-4} & (Q) & Metallic precipitate is obtained \\
\hline
(c) \textbf{STEP-7} & (R) & Dark blue precipitate is formed \\
\hline
(d) \textbf{STEP-8} & (S) & All four metal ions are present in precipitate, No ion in solution \\
\hline
\end{tabular}}\par
The correct code is :-\par
\begin{qoptions}
\opt{A}{a$\rightarrow$S , b$\rightarrow$Q , c$\rightarrow$P , d$\rightarrow$R}
\opt{B}{a$\rightarrow$S , b$\rightarrow$R , c$\rightarrow$P , d$\rightarrow$Q}
\opt{C}{a$\rightarrow$S , b$\rightarrow$P , c$\rightarrow$Q , d$\rightarrow$R}
\opt{D}{a$\rightarrow$P , b$\rightarrow$S , c$\rightarrow$R , d$\rightarrow$Q}
\end{qoptions}
\end{bankq}
\begin{bankq}[id=SA15-P1-CHEM-S2-Q12, source={Score Advanced - Test 15 (18 May 2023)}, paper=1, qno=12, type=multi, subtopic={Anion tests}, answer={A,B,C}, answer_src=solved, working={A, B and C are correct}, difficulty=10, diff_parts={K3 S3 I2 C2 F1}, diff_note={concepts: identification of barium nitrite and ammonium sulphate, anion tests, A is Ba(NO2)2 | steps: A is Ba(NO2)2; B is (NH4)2SO4; D is BaO; E is BaCO3; iodine dissolves as I3-}, also_in={Score Advanced - Test 32 (08 May 2024) P1 Q12}]
A (is a compound of Ba) and B (is a sulphate) are white crystalline substances, highly soluble in water and decompose at higher temperatures. If an aqueous solution of A (slightly basic) is added to an aqueous solution of B (slightly acidic) a white precipitate C forms. The filtrate is neutral and gives a brown colour solution with an excess of acidified KI solution. When boiled, the filtrate evaporates without the appearance of any residue. White solid D can be prepared by the heating of A in the absence of air. The reaction of D with water gives a colourless solution. This solution, if kept in an open container, slowly precipitates a white solid E and leaves water. Upon prolonged exposure to moist air at room temperature, solid D is transformed into E as well. E can be converted back into D but decomposition above 1400\textdegree C is required for this purpose. The reaction of B and D in water forms the precipitate C and is accompanied by a characteristic odour due to formation of a gas. A gives white precipitate with AgNO$_3$. Select the correct option(s):
\begin{qoptions}
\opt{A}{Compound A is on heating gives neutral, acidic and basic oxide}
\opt{B}{Compound B is (NH$_4$)$_2$SO$_4$.}
\opt{C}{Compound E is BaCO$_3$.}
\opt{D}{I$_2$ will precipitate as brown coloured.}
\end{qoptions}
\end{bankq}
\begin{bankpassage}[id=SA16-P2-CHEM-S3-G5-6, source={Score Advanced - Test 16 (20 May 2023)}, questions={5-6}]
Consider following reagents with given codes :\par
\textbf{Code 1 :} dil.HCl\par
\textbf{Code 2 :} NH$_3$ (aq)\par
\textbf{Code 3 :} KI (aq)\par
\textbf{Code 4 :} KCN (aq)\par
\textbf{Code 5 :} Na$_2$S$_2$O$_3$ (aq)\par
\textbf{Code 6 :} NaOH (aq)\par
Let the sum of codes of reagents in which AgCl is soluble be $x$.\par
Let the sum of codes of reagents in which Zn(OH)$_2$ is soluble be $y$\par
\end{bankpassage}
\begin{bankq}[id=SA16-P2-CHEM-S3-Q05, source={Score Advanced - Test 16 (20 May 2023)}, paper=2, qno=5, type=numerical, subtopic={Cation tests}, answer={5.50}, answer_src=key, difficulty=7, diff_parts={K3 S1 I2 C1 F1}, diff_note={concepts: sum of codes of reagents dissolving silver chloride, cation tests, complex formation with NH3 thiosulphate cyanide | steps: complex formation with NH3 thiosulphate cyanide}, group=SA16-P2-CHEM-S3-G5-6]
Find $\dfrac{x}{2}$
\end{bankq}
\begin{bankq}[id=SA16-P2-CHEM-S3-Q06, source={Score Advanced - Test 16 (20 May 2023)}, paper=2, qno=6, type=numerical, subtopic={Cation tests}, answer={3.25}, answer_src=key, difficulty=7, diff_parts={K3 S1 I2 C1 F1}, diff_note={concepts: sum of codes of reagents dissolving zinc hydroxide, cation tests, amphoteric | steps: amphoteric; ammine complex}, group=SA16-P2-CHEM-S3-G5-6]
Find $\dfrac{y}{4}$
\end{bankq}
\begin{bankq}[id=SA17-P1-CHEM-S1-Q01, source={Score Advanced - Test 17 (25 May 2023)}, paper=1, qno=1, type=multi, subtopic={Cation tests}, answer={A,B,C,D}, answer_src=key, difficulty=7, diff_parts={K3 S1 I2 C1 F1}, diff_note={concepts: precipitate formation of metal salts, cation tests, check each reaction for a precipitate | steps: check each reaction for a precipitate}]
How many of the following reactions will not result in precipitate formation ?\par
(ignore balancing)
\begin{qoptions}
\opt{A}{FeCl$_3$(aq) + Na$_2$SO$_4$(aq) $\rightarrow$}
\opt{B}{Hg(NO$_3$)$_3$(aq) + exNaCN(aq) $\rightarrow$}
\opt{C}{NiCl$_2$(aq) + ex NH$_3$(aq) $\rightarrow$}
\opt{D}{ZnCl$_2$(aq) + H$_2$S(aq) $\xrightarrow{\text{H}^+}$}
\end{qoptions}
\end{bankq}
\begin{bankq}[id=SA17-P1-CHEM-S3-Q02, source={Score Advanced - Test 17 (25 May 2023)}, paper=1, qno=2, type=single, subtopic={Cation tests}, answer={1.50}, answer_src=key, difficulty=6, diff_parts={K3 S1 I2 C1 F0}, diff_note={concepts: precipitates soluble in cyanide and in dilute acid, cation tests, solubility of sulphides and cyanides | steps: solubility of sulphides and cyanides}, also_in={Score Advanced - Test 35 (15 May 2024) P1 Q2}]
Consider following precipitates :\par
CuS, ZnS, HgS, Ni(CN)$_2$, AgCN, AgCl, CuCN, Pb(CN)$_2$, FeS, MnS\par
Let,\par
$x=$ number of precipitates soluble in aq.KCN\par
$y=$ number of precipitates soluble in dil.HCl\par
Find $\dfrac{|x-y|}{2}$
\end{bankq}
\begin{bankq}[id=SA17-P2-CHEM-S3-Q01, source={Score Advanced - Test 17 (25 May 2023)}, paper=2, qno=1, type=numerical, subtopic={Cation tests}, answer={3.00}, answer_src=key, difficulty=7, diff_parts={K3 S1 I2 C1 F1}, diff_note={concepts: numbers of yellow and brown compounds, cation tests, colour of each compound | steps: colour of each compound}, also_in={Score Advanced - Test 35 (15 May 2024) P2 Q1}]
Consider following compounds :\par
MnO$_2$, HgI$_2$, Cu(CN)$_2$, Cu$_2$[Fe(CN)$_6$], Fe(CN)$_3$, HgSO$_4$.2HgO, Cu$_2$O, MnS, [Ni(CN)$_4$]$^{2-}$\par
Let,\par
$x$ = number of yellow coloured compounds\par
$y$ = number of brown coloured compounds\par
Find $\frac{xy}{2}$
\end{bankq}
\begin{bankq}[id=SA19-P1-CHEM-S1-Q02, source={Score Advanced - Test 19 (04 Apr 2024)}, paper=1, qno=2, type=single, subtopic={Cation tests}, answer={D}, answer_src=key, difficulty=7, diff_parts={K3 S1 I2 C2 F0}, diff_note={concepts: group analysis statements, cation tests, group III precipitation with ammonium chloride | steps: group III precipitation with ammonium chloride; sulphate test}]
Identify the CORRECT statements.\par
(I) NH$_4$Cl is added along with NH$_4$OH so that only III$^{rd}$ group cations can be precipitated as their hydroxides.\par
(II) In III$^{rd}$ group in place of NH$_4$Cl, (NH$_4$)$_2$SO$_4$ cannot be used as V$^{th}$ group radical(s) will also be precipitated as MSO$_4$ along with Al$^{3+}$, Fe$^{3+}$ \& Cr$^{3+}$.\par
(III) Aqueous solution of ammonium sulphate produces white ppt. with BaCl$_2$ solution which is insoluble in dil.HCl.
\begin{qoptions}
\opt{A}{I and II only}
\opt{B}{II and III only}
\opt{C}{I and III only}
\opt{D}{I, II and III}
\end{qoptions}
\end{bankq}
\begin{bankpassage}[id=SA21-P1-CHEM-S2-G7-8, source={Score Advanced - Test 21 (08 Apr 2024)}, questions={7-8}]
A sample (X) of compound(s) contains any one or two of the following AgNO$_3$, ZnSO$_4$,CuSO$_4$, MgCl$_2$, CaCO$_3$, BaCl$_2$, Na$_2$SO$_4$, NaCl, NaOH\par
When (X) is dissolved in water, a clear solution is obtained. This solution is divided into two parts I and II. Part I turns pink on addition of a few drops of phenolphthalein. When dilute hydrochloric acid is added slowly to part II, a white precipitate (Y) is obtained which dissolves on further addition of dilute hydrochloric acid
\end{bankpassage}
\begin{bankq}[id=SA21-P1-CHEM-S2-Q07, source={Score Advanced - Test 21 (08 Apr 2024)}, paper=1, qno=7, type=single, subtopic={Cation tests}, answer={C}, answer_src=key, difficulty=7, diff_parts={K3 S1 I2 C2 F0}, diff_note={concepts: species present in the clear solution of the sample, cation tests, alkaline solution gives zincate | steps: alkaline solution gives zincate}, group=SA21-P1-CHEM-S2-G7-8]
The clear solution obtained when (X) is dissolved in water, contains
\begin{qoptions}
\opt{A}{ZnCl$_4^{2-}$}
\opt{B}{AgCl$_2^{-}$}
\opt{C}{Zn(OH)$_4^{2-}$}
\opt{D}{CuCl$_4^{2-}$}
\end{qoptions}
\end{bankq}
\begin{bankq}[id=SA21-P1-CHEM-S2-Q08, source={Score Advanced - Test 21 (08 Apr 2024)}, paper=1, qno=8, type=single, subtopic={Cation tests}, answer={D}, answer_src=key, difficulty=7, diff_parts={K3 S1 I2 C2 F0}, diff_note={concepts: identify the two compounds from tests, cation tests, phenolphthalein | steps: phenolphthalein; white precipitate soluble in acid}, group=SA21-P1-CHEM-S2-G7-8]
The sample (X) contains
\begin{qoptions}
\opt{A}{AgNO$_3$ \& NaCl}
\opt{B}{ZnSO$_4$ \& BaCl$_2$}
\opt{C}{HgCl$_2$ \& NaOH}
\opt{D}{ZnSO$_4$ \& NaOH}
\end{qoptions}
\end{bankq}
\begin{bankq}[id=SA22-P1-CHEM-S3-Q02, source={Score Advanced - Test 22 (10 Apr 2024)}, paper=1, qno=2, type=numerical, subtopic={Anion tests}, answer={7}, answer_src=key, difficulty=9, diff_parts={K3 S2 I2 C2 F1}, diff_note={concepts: identify gases from heating a metal sulphate, anion tests, SO2 reduces dichromate | steps: SO2 reduces dichromate; SO3 trimerises; count}]
A sulphate of a metal (M) on heating evolves two gases (B) and (C) and an oxide (D). Gas (B), turns acidified dichromate paper green while gas (C) forms a trimer.\par
Compound (D) forms a Lewis acid (E) with HCl (g). If the number of S $-$ O $-$ S bonds in trimer of gas (C) be `n' and the number of equivalent M $-$ Cl bonds in one molecule of (E) in vapour phase be `m' then (m+n) is
\end{bankq}
\begin{bankpassage}[id=SA23-P2-CHEM-S3-G1-2, source={Score Advanced - Test 23 (14 Apr 2024)}, questions={1-2}]
An inorganic compound (A) reacts with dilute H$_2$SO$_4$ to give a colourless gas (B) which smells like rotten egg.\par
Black precipitate of (X) $\xleftarrow{\text{AgNO}_3}$ B $\xrightarrow{\text{Alk.Na}_2[\text{Fe(CN)}_5\text{NO}]}$ Violet coloured compound (Y)
\end{bankpassage}
\begin{bankq}[id=SA23-P2-CHEM-S3-Q01, source={Score Advanced - Test 23 (14 Apr 2024)}, paper=2, qno=1, type=numerical, subtopic={Anion tests}, answer={0.00}, answer_src=key, difficulty=8, diff_parts={K3 S1 I2 C2 F1}, diff_note={concepts: change in oxidation state of iron in nitroprusside test, anion tests, sulphide attacks NO | steps: sulphide attacks NO; iron stays +2}, group=SA23-P2-CHEM-S3-G1-2]
Change in oxidation state of Fe, When Na$_2[\text{Fe(CN)}_5\text{NO}]$ reacts with B to give (Y).
\end{bankq}
\begin{bankq}[id=SA23-P2-CHEM-S3-Q02, source={Score Advanced - Test 23 (14 Apr 2024)}, paper=2, qno=2, type=numerical, subtopic={Anion tests}, answer={3.33}, answer_src=key, difficulty=8, diff_parts={K3 S1 I2 C2 F1}, diff_note={concepts: moles of dichromate to oxidise H2S, equivalents, moles of H2S from the mass | steps: moles of H2S from the mass; electrons per mole}, group=SA23-P2-CHEM-S3-G1-2]
Number of moles of acidified K$_2$Cr$_2$O$_7$ required to oxidise 340 gram gas B.\par
(Given atomic mass : H = 1, O = 16, C = 12, F = 19, N = 14, S = 32, P = 31, Cl = 35.5, Br = 80, I = 127, K = 39, Na = 23, Pb = 207, Al = 27, Si = 28, Mn = 55, Cr = 52)
\end{bankq}
\begin{bankq}[id=SA25-P1-CHEM-S4-Q04, source={Score Advanced - Test 25 (21 Apr 2024)}, paper=1, qno=4, type=numerical, subtopic={Cation tests}, answer={331}, answer_src=key, difficulty=8, diff_parts={K3 S1 I2 C2 F1}, diff_note={concepts: molecular weight of a white precipitate in a reaction scheme, cation tests, identify the cation reagents | steps: identify the cation reagents; precipitate}, also_in={Score Advanced - Test 26 (21 Apr 2024) P1 Q4}]
Consider the following reactions and write molecular weight of \textbf{white precipitate (Y)}.\par
(Given atomic mass : H = 1, O = 16, C = 12, F = 19, N = 14, S = 32, P = 31, Cl = 35.5, Br = 80, I = 127, K = 39, Na = 23, Pb = 207, Al = 27, Si = 28, Mn = 55, Cr = 52)\par
\bfig{229.5}{SA25_p1_chem_s4_q04_1.png}
Colored gas [Q] is identified by its staining of starch paper orange red\par
Phenol + [Q] in water $\rightarrow$ \textbf{White precipitate (Y)} + other product
\end{bankq}
\begin{bankq}[id=SA27-P1-CHEM-S1-Q04, source={Score Advanced - Test 27 (24 Apr 2024)}, paper=1, qno=4, type=single, subtopic={Anion tests}, answer={B}, answer_src=key, difficulty=7, diff_parts={K3 S1 I2 C2 F0}, diff_note={concepts: salt giving a red colour after cyanide treatment, anion tests, thiocyanate formation from sulphur | steps: thiocyanate formation from sulphur}]
To an alkaline salt solution KCN added. Now the solution acidified and FeCl$_3$ added, then red colouration observed. The salt may posses
\begin{qoptions}
\opt{A}{SO$_3^{2-}$}
\opt{B}{S$_2$O$_3^{2-}$}
\opt{C}{NO$_2^-$}
\opt{D}{SO$_4^{2-}$}
\end{qoptions}
\end{bankq}
\begin{bankq}[id=SA27-P2-CHEM-S1-Q04, source={Score Advanced - Test 27 (24 Apr 2024)}, paper=2, qno=4, type=multi, subtopic={Cation tests}, answer={B,C,D}, answer_src=solved, working={KSCN gives red colour with ferric ions, not ferrous. Ferrous ions give red complexes with dimethylglyoxime in ammonia, o-phenanthroline (ferroin) and 2,2'-bipyridyl. B, C and D.}, difficulty=8, diff_parts={K3 S1 I2 C2 F1}, diff_note={concepts: reagents giving a red colour with ferrous ions in the presence of ferric ions, cation tests, ferrous specific reagents | steps: ferrous specific reagents}]
With which of the following reagents, ferrous ions give red colouration in the presence of ferric ions
\begin{qoptions}
\opt{A}{KSCN}
\opt{B}{NH$_3$ + dmg}
\opt{C}{O-Phenanthroline}
\opt{D}{2, 2$'$-Bipyridyl}
\end{qoptions}
\end{bankq}
\begin{bankq}[id=SA30-P2-CHEM-S3-Q01, source={Score Advanced - Test 30 (04 May 2024)}, paper=2, qno=1, type=numerical, subtopic={Cation tests}, answer={4}, answer_src=key, difficulty=8, diff_parts={K3 S1 I2 C2 F1}, diff_note={concepts: standard problem in this chapter, cation tests, apply the relevant results | steps: apply the relevant results; compute}, also_in={Score Advanced - Test 31 (04 May 2024) P2 Q1}]
How many of the following cations gives ppt with H$_2$S in acidic solution :-\par
Zn$^{2+}$ , Cu$^{2+}$ , Ni$^{2+}$ , Sr$^{2+}$ , Al$^{3+}$ , Bi$^{3+}$ , Sn$^{2+}$ , Pb$^{2+}$
\end{bankq}
\begin{bankq}[id=SA32-P2-CHEM-S2-Q05, source={Score Advanced - Test 32 (08 May 2024)}, paper=2, qno=5, type=multi, subtopic={Cation tests}, answer={A,B,C}, answer_src=key, difficulty=8, diff_parts={K3 S1 I2 C2 F1}, diff_note={concepts: standard problem in this chapter, cation tests, apply the relevant results | steps: apply the relevant results; compute}]
Cation(s) which form an ammine complex ion with excess of NH$_3$ is/are
\begin{qoptions}
\opt{A}{Ag$^+$}
\opt{B}{Cu$^{+2}$}
\opt{C}{Co$^{+2}$}
\opt{D}{Al$^{+3}$}
\end{qoptions}
\end{bankq}
\begin{bankq}[id=SA33-P2-CHEM-S1-Q03, source={Score Advanced - Test 33 (12 May 2024)}, paper=2, qno=3, type=single, subtopic={Cation tests}, answer={B}, answer_src=key, difficulty=4, diff_parts={K2 S1 I1 C1 F0}, diff_note={concepts: standard problem in this chapter, cation tests | steps: apply the relevant results; compute}, also_in={Score Advanced - Test 34 (12 May 2024) P2 Q3}]
CoS (Black) $\xrightarrow{\text{aque--regia}}$ solution $\xrightarrow[\mathrm{Br_2\ water}]{\mathrm{exc.NaHCO_3}}$ coloured complex (P)\par
What is co-ordination no. of central metal in complex (P)
\begin{qoptions}
\opt{A}{5}
\opt{B}{6}
\opt{C}{3}
\opt{D}{4}
\end{qoptions}
\end{bankq}
\begin{bankq}[id=SA33-P2-CHEM-S3-Q05, source={Score Advanced - Test 33 (12 May 2024)}, paper=2, qno=5, type=numerical, subtopic={Cation tests}, answer={5}, answer_src=key, difficulty=8, diff_parts={K3 S1 I2 C2 F1}, diff_note={concepts: standard problem in this chapter, cation tests, apply the relevant results | steps: apply the relevant results; compute}, also_in={Score Advanced - Test 34 (12 May 2024) P2 Q5}]
A solution contains equimolar amount of following ions : Zn$^{+2}$, Fe$^{+3}$, Ag$^+$, Cu$^{+2}$, Ni$^{+2}$, Pb$^{+2}$, Cd$^{+2}$, Hg$^{+2}$, Cr$^{+3}$. How many of cations form soluble complex with excess of aqueous NH$_3$?
\end{bankq}
\begin{bankq}[id=SA34-P1-CHEM-S3-Q04, source={Score Advanced - Test 34 (12 May 2024)}, paper=1, qno=4, type=single, subtopic={Cation tests}, answer={14}, answer_src=key, difficulty=7, diff_parts={K3 S1 I2 C2 F0}, diff_note={concepts: standard problem in this chapter, cation tests, apply the relevant results | steps: apply the relevant results; compute}, also_in={Score Advanced - Test 36 (19 May 2024) P1 Q4}]
A mixture of Potash alum, Chrome alum and Mohr's salt is taken in water and H$_2$O$_2$ is added in strongly alkaline medium and finally filtered.\par
If atomic mass/molecular mass of cation in residue is ``a'' and the number of type of ions from the salt mixture passed into the filtrate is `b'. Then find the value of a/b.\par
(Atomic mass : Cr = 52 , Fe = 56 , S = 32 , O = 16 , N = 14 , C = 12 , B = 11 , H = 1 , K$^+$ = 39 , Na$^+$ = 23 , Ni = 58 , Cu = 29 , Zn = 65 , Cl = 35.5 , Br = 80 , F = 19)
\end{bankq}
\begin{bankq}[id=SA35-P1-CHEM-S1-Q01, source={Score Advanced - Test 35 (15 May 2024)}, paper=1, qno=1, type=multi, subtopic={Cation tests}, answer={A,B,C,D}, answer_src=key, difficulty=7, diff_parts={K3 S1 I2 C1 F1}, diff_note={concepts: precipitate formation of metal salts, cation tests, check each reaction for a precipitate | steps: check each reaction for a precipitate}]
How many of the following reactions will not result in precipitate formation ?\par
(ignore balancing)
\begin{qoptions}
\opt{A}{FeCl$_3$(aq) + Na$_2$SO$_4$(aq) $\to$}
\opt{B}{Hg(NO$_3$)$_2$(aq) + exNaCN(aq) $\to$}
\opt{C}{NiCl$_2$(aq) + ex NH$_3$(aq) $\to$}
\opt{D}{ZnCl$_2$(aq) + H$_2$S(aq) $\xrightarrow{H^+}$}
\end{qoptions}
\end{bankq}
\begin{bankq}[id=SA36-P2-CHEM-S4-Q09, source={Score Advanced - Test 36 (19 May 2024)}, paper=2, qno=9, type=numerical, subtopic={Cation tests}, answer={3}, answer_src=key, difficulty=8, diff_parts={K3 S1 I2 C2 F1}, diff_note={concepts: standard problem in this chapter, cation tests, apply the relevant results | steps: apply the relevant results; compute}, also_in={Score Advanced - Test 37 (19 May 2024) P2 Q9}]
Consider the following precipitate/s\par
PbS , CuS , HgS , CdS , As$_2$S$_3$, As$_2$S$_5$, SnS, SnS$_2$, Sb$_2$S$_3$, Sb$_2$S$_5$\par
x = Total number of yellow ppt/s which soluble in yellow ammonium sulphide\par
y = Total number of black ppt/s which insoluble in hot \& dil. HNO$_3$\par
Find the value of $\dfrac{\mathrm x}{\mathrm y}$.
\end{bankq}
\begin{bankpassage}[id=SA38-P2-CHEM-S3-G3-4, source={Score Advanced - Test 38 (13 Apr 2025)}, questions={3-4}]
When potassium iodide is added to an aqueous solution of potassium ferricyanide, a reversible reaction is observed in which a complex \textbf P is formed. In a strong acidic medium, the equilibrium shifts completely towards \textbf P. Addition of ferric chloride to \textbf P in a slightly acidic medium results in insoluble complex \textbf Q.
\end{bankpassage}
\begin{bankq}[id=SA38-P2-CHEM-S3-Q03, source={Score Advanced - Test 38 (13 Apr 2025)}, paper=2, qno=3, type=numerical, subtopic={Anion tests}, answer={9}, answer_src=key, difficulty=8, diff_parts={K3 S1 I2 C2 F1}, diff_note={concepts: standard problem in this chapter, anion tests, apply the relevant results | steps: apply the relevant results; compute}, group=SA38-P2-CHEM-S3-G3-4]
If 3 moles of potassium iodide are taken then X moles of P and Y moles of iodine is formed. Find the value of (X+4Y).
\end{bankq}
\begin{bankq}[id=SA38-P2-CHEM-S3-Q04, source={Score Advanced - Test 38 (13 Apr 2025)}, paper=2, qno=4, type=numerical, subtopic={Anion tests}, answer={7}, answer_src=solved, working={2 [Fe(CN)6]3- + 2 I- gives 2 [Fe(CN)6]4- + I2; in strong acid P is H4[Fe(CN)6]. Ferric chloride gives insoluble Prussian blue Q = Fe4[Fe(CN)6]3 which contains 4 + 3 = 7 iron atoms.}, difficulty=8, diff_parts={K3 S1 I2 C2 F1}, diff_note={concepts: standard problem in this chapter, transition elements, apply the relevant results | steps: apply the relevant results; compute}, group=SA38-P2-CHEM-S3-G3-4]
The number of iron present in the molecular formula of \textbf Q is \underline{\hspace{1.5cm}}.
\end{bankq}
\begin{bankq}[id=SA39-P2-CHEM-S1-Q05, source={Score Advanced - Test 39 (20 Apr 2025)}, paper=2, qno=5, type=multi, subtopic={Anion tests}, answer={Bonus}, answer_src=solved, working={Printed key for this paper is incomplete and the figure/data needed to solve it was not recoverable; banked as Bonus.}, difficulty=5, diff_parts={K2 S1 I1 C1 F1}, diff_note={concepts: standard problem in this chapter, anion tests | steps: apply the relevant results; compute}]
A metal cation M$^{+2}$ gives black precipitate with acidified H$_2$S which dissolves in hot conc. HNO$_3$. Solution of M$^{+2}$ gives brown precipitate on treatment with Potassium ferrocyanide. Select the correct statement(s) :
\begin{qoptions}
\opt{A}{M$^{+2}$(aq.) on reaction with exc. KCN produces a gas.}
\opt{B}{Metal `M' in its extraction process from its most common sulphide ore is obtained from auto-reaction process.}
\opt{C}{E$^\circ_{\mathrm M^{+2}/\mathrm M}$ is positive.}
\opt{D}{M$^{2+}_{\mathrm{(aq.)}}$ produce dark blue solution with excess Cl$^-$}
\end{qoptions}
\end{bankq}
\begin{bankq}[id=SA40-P2-CHEM-S4-Q04, source={Score Advanced - Test 40 (27 Apr 2025)}, paper=2, qno=4, type=single, subtopic={Anion tests}, answer={4}, answer_src=key, difficulty=7, diff_parts={K3 S1 I2 C2 F0}, diff_note={concepts: standard problem in this chapter, anion tests, apply the relevant results | steps: apply the relevant results; compute}]
Find the number of combination(s) of reagents in which pink coloured solution or precipitate or product(s) is/are formed.\par
(i) MnCl$_2$ + H$_2$S(g) + HCl\qquad (ii) Mn$^{2+}$ + S$_2$O$_8^{2-}$\par
(iii) Mn$^{2+}$ + NaBiO$_3$ + HNO$_3$\qquad (iv) Mn$^{2+}$ + PbO$_2$ + HNO$_3$\par
(v) K$_2$MnO$_4$ + Cl$_2$\qquad (vi) ZnCl$_2$ + H$_2$S + NH$_4$OH(aq.)
\end{bankq}
\begin{bankq}[id=SA41-P2-CHEM-S4-Q07, source={Score Advanced - Test 41}, paper=2, qno=7, type=numerical, subtopic={Anion tests}, answer={5}, answer_src=key, difficulty=8, diff_parts={K3 S1 I2 C2 F1}, diff_note={concepts: standard problem in this chapter, anion tests, apply the relevant results | steps: apply the relevant results; compute}]
Number of correct statements among following is.\par
(I) Barium salt imparts apple green colour to oxidising flame\par
(II) In Mac-Arthur Forest cyanide process Zn act as reducing agent\par
(III) K$_2$CO$_3$ cannot be prepared by Solvay soda process\par
(IV) Fe(CO)$_5$ is diamagnetic and having trigonal bipyramidal geometry\par
(V) [Co(H$_2$O)$_3$F$_3$] is high spin complex\par
(VI) EA$_1$ : F $>$ Cl $>$ Br $>$ I\par
(VII) Bond energy : Cl$_2$ $>$ F$_2$ $>$ Br$_2$ $>$ I$_2$
\end{bankq}

\section*{SRGBT papers}
\begin{bankq}[id=SRGBT13-P1-CHEM-S1-Q29, source={SRGBT - Test 13 (17 Oct 2021)}, paper=1, qno=29, type=multi, subtopic={Anion tests}, answer={B}, answer_src=key, difficulty=5, diff_parts={K2 S1 I1 C1 F1}, diff_note={concepts: standard problem in this chapter, anion tests | steps: apply the relevant results; compute}]
Select correct ``reaction and observation'' combination(s).\par
\begin{tabular}{@{}l p{3.4cm} p{3.4cm}@{}}
 & \underline{\textbf{Reactions}} & \underline{\textbf{Observations}}\\
(i) & NaCl + NH$_4$HCO$_3$ $\longrightarrow$ Product(s) & Polyatomic anion in product species has sp$^3$ hybrid central atom\\
(ii) & Na$_2$CO$_3$ + Ca(OH)$_2$ $\longrightarrow$ Product(s) & One of the products is insoluble precipitate\\
(iii) & KOH + O$_3$ $\longrightarrow$ Product(s) & Products include - a neutral oxide, two paramagnetic species\\
\end{tabular}\par
\begin{qoptions}
\opt{A}{(i) , (ii) only}
\opt{B}{(ii) , (iii) only}
\opt{C}{(ii) only}
\opt{D}{(i) , (ii) \& (iii)}
\end{qoptions}
\end{bankq}

\section*{Enthusiast Pre-Score papers}
\begin{bankq}[id=EPS02-P1-CHEM-S1-Q01, source={Enthusiast Pre-Score - Test 02 (11 Nov 2020)}, paper=1, qno=1, type=single, subtopic={Cation tests}, answer={D}, answer_src=key, difficulty=5, diff_parts={K2 S2 I1 C1 F0}, diff_note={concepts: analysis of anions and cations from precipitates, cation tests | steps: identify the anion by the precipitate; identify cations by flame; check each}]
(1) Bottle-1 (Colourless solution) + Pb(OAc)$_2$ solution $\longrightarrow$ white ppt. which dissolves only on strong heating ($\sim$85-90$^\circ$C)\par
(2) Bottle-2 (Coloured solution) + Pb(OAc)$_2$ solution $\longrightarrow$ yellow ppt. insoluble in acidic acid.\par
(3) Solid salts in bottle-1 and bottle-2 are mixed together and taking that mixture, flame test is performed. Voluminous yellow flame is observed with naked eye and Crimson red flame is observed through double blue glass.\par
Which of the following statement is \textbf{INCORRECT}?
\begin{qoptions}
\opt{A}{The anion present in bottle-2 is $\mathrm{CrO_4^{2-}}$}
\opt{B}{The cation present in bottle-1 may be Na$^+$}
\opt{C}{The cation present in bottle-2 may be Na$^+$}
\opt{D}{The anion present in bottle-1 is Cl$^-$}
\end{qoptions}
\end{bankq}
\begin{bankq}[id=EPS02-P1-CHEM-S1-Q02, source={Enthusiast Pre-Score - Test 02 (11 Nov 2020)}, paper=1, qno=2, type=single, subtopic={Cation tests}, answer={D}, answer_src=key, difficulty=2, diff_parts={K1 S1 I1 C0 F0}, diff_note={concepts: colour of group III precipitates | steps: identify the precipitate; oxidation state}]
In the group analysis of basic radicals, the colour of the ppt. obtained in group-III test is reddish brown. The oxidation state of cation present in the original salt may be.
\begin{qoptions}
\opt{A}{+2}
\opt{B}{+3}
\opt{C}{+2 and +3}
\opt{D}{All of these}
\end{qoptions}
\end{bankq}
\begin{bankq}[id=EPS02-P1-CHEM-S2-Q07, source={Enthusiast Pre-Score - Test 02 (11 Nov 2020)}, paper=1, qno=7, type=multi, subtopic={Cation tests}, answer={B,C,D}, answer_src=key, difficulty=4, diff_parts={K1 S2 I1 C0 F1}, diff_note={concepts: properties of the iodide of Millon base | steps: structure; formation; colour}]
Which of the following statements is/are \textbf{CORRECT} for iodide of Millon's base?
\begin{qoptions}
\opt{A}{The formation reaction of it is a redox reaction}
\opt{B}{The number atoms per molecule is 7}
\opt{C}{The colour of it is brown}
\opt{D}{The test associated to it is used to detect NH$_3$ gas}
\end{qoptions}
\end{bankq}
\begin{bankq}[id=EPS02-P1-CHEM-S2-Q08, source={Enthusiast Pre-Score - Test 02 (11 Nov 2020)}, paper=1, qno=8, type=multi, subtopic={Cation tests}, answer={B,D}, answer_src=key, difficulty=4, diff_parts={K1 S2 I1 C0 F1}, diff_note={concepts: cations giving yellow sulphides soluble in yellow ammonium sulphide | steps: group II B sulphides; colour; solubility}]
Salt solution `A' is treated with H$_2$S in presence of dil. HCl and resulting a yellow ppt. which dissolves in YAS. The cation present in the salt may be.
\begin{qoptions}
\opt{A}{Sn$^{2+}$}
\opt{B}{Sn$^{4+}$}
\opt{C}{Cd$^{2+}$}
\opt{D}{As$^{+5}$}
\end{qoptions}
\end{bankq}
\begin{bankq}[id=EPS02-P1-CHEM-S3-Q01, source={Enthusiast Pre-Score - Test 02 (11 Nov 2020)}, paper=1, qno=1, type=numerical, subtopic={Anion tests}, answer={0.665}, answer_src=key, difficulty=5, diff_parts={K1 S2 I1 C1 F1}, diff_note={concepts: statements on qualitative tests | steps: layer test; flame and group tests; distinguishing reagents; count}]
Find the ratio of number of incorrect to correct statements from the following.\par
(i) In the layer test, the organic layer is used to intensify the observation.\par
(ii) Layer test is the confimation test for Br$^-$ and Cl$^-$.\par
(iii) K$_4$[Fe(CN)$_6$] can be used to distinguish between Ca$^{2+}$ and Sr$^{+2}$\par
(iv) The addition of Na$_2$C$_2$O$_4$ solution followed by acetic acid can distinguish between Ca$^{2+}$ and Ba$^{2+}$\par
(v) K$^+$ and NH$_4^+$ can be distinguished by the addition of H$_2$[PtCl$_6$].
\end{bankq}
\begin{bankq}[id=EPS02-P2-CHEM-S1-Q01, source={Enthusiast Pre-Score - Test 02 (11 Nov 2020)}, paper=2, qno=1, type=multi, subtopic={Cation tests}, answer={A,D}, answer_src=key, difficulty=2, diff_parts={K1 S1 I0 C0 F1}, diff_note={concepts: precipitates with potassium ferrocyanide | steps: colours of the metal ferrocyanides}]
Potassium ferrocyanide is an important laboratory reagent used to detect many cations. Select the metal ions that can give almost white ppt. with it.
\begin{qoptions}
\opt{A}{Zn$^{2+}$}
\opt{B}{Fe$^{3+}$}
\opt{C}{Cu$^{2+}$}
\opt{D}{Fe$^{2+}$}
\end{qoptions}
\end{bankq}
\begin{bankq}[id=EPS02-P2-CHEM-S1-Q02, source={Enthusiast Pre-Score - Test 02 (11 Nov 2020)}, paper=2, qno=2, type=multi, subtopic={Cation tests}, answer={A,C,D}, answer_src=key, difficulty=3, diff_parts={K1 S1 I1 C0 F1}, diff_note={concepts: zinc and aluminium analysis | steps: reaction with ammonia and sodium hydroxide; sulphide precipitation}]
Which of the following are ``True'' statements?
\begin{qoptions}
\opt{A}{Zn and Al salts can be distinguished by using NH$_4$OH.}
\opt{B}{Zn and Al salts can be distinguished by using NaOH}
\opt{C}{ZnS is precipitated by H$_2$S in ammonical solution of zinc salt but not in the acidic solution.}
\opt{D}{Al(OH)$_3$ is soluble in NaOH as well as HCl}
\end{qoptions}
\end{bankq}
\begin{bankq}[id=EPS02-P2-CHEM-S2-Q09, source={Enthusiast Pre-Score - Test 02 (11 Nov 2020)}, paper=2, qno=9, type=single, subtopic={Anion tests}, answer={C}, answer_src=key, difficulty=2, diff_parts={K1 S1 I1 C0 F0}, diff_note={concepts: anion tests table | steps: tests for each anion; match}]
\textbf{Answer the following by appropriately matching the lists based on the information given in the paragraph.}\par
\smallskip
\noindent{\footnotesize\setlength{\tabcolsep}{2pt}\begin{tabular}{|L{.8cm}|L{2.0cm}|L{.8cm}|L{3.3cm}|}\hline
\multicolumn{2}{|c|}{List-I: Radicals} & \multicolumn{2}{c|}{\parbox{3.9cm}{\centering List-II: Correct statement regarding them}}\\\hline
(I) & Br$^-$ & (P) & Can be confirmed by Brown ring test\\\hline
(II) & I$^-$ & (Q) & Can be confirmed by Layer test\\\hline
(III) & S$^{2-}$ & (R) & Can be confirmed by methylene blue test\\\hline
(IV) & NO$_3^-$ & (S) & On reaction with conc. H$_2$SO$_4$, oxidation takes place\\\hline
 & & (T) & All salts are soluble\\\hline
 & & (U) & On reaction with dilute H$_2$SO$_4$ produces colourless gas\\\hline
\end{tabular}}\par
Which of the following option has the \textbf{CORRECT} combination considering List-I \& List-II ?
\begin{qoptions}
\opt{A}{I, P}
\opt{B}{II, T}
\opt{C}{III, R}
\opt{D}{IV, S}
\end{qoptions}
\end{bankq}
\begin{bankq}[id=EPS02-P2-CHEM-S2-Q10, source={Enthusiast Pre-Score - Test 02 (11 Nov 2020)}, paper=2, qno=10, type=single, subtopic={Anion tests}, answer={B}, answer_src=key, difficulty=2, diff_parts={K1 S1 I1 C0 F0}, diff_note={concepts: anion tests table incorrect match | steps: tests for each anion; find the wrong match}]
\textbf{Answer the following by appropriately matching the lists based on the information given in the paragraph.}\par
\smallskip
\noindent{\footnotesize\setlength{\tabcolsep}{2pt}\begin{tabular}{|L{.8cm}|L{2.0cm}|L{.8cm}|L{3.3cm}|}\hline
\multicolumn{2}{|c|}{List-I: Radicals} & \multicolumn{2}{c|}{\parbox{3.9cm}{\centering List-II: Correct statement regarding them}}\\\hline
(I) & Br$^-$ & (P) & Can be confirmed by Brown ring test\\\hline
(II) & I$^-$ & (Q) & Can be confirmed by Layer test\\\hline
(III) & S$^{2-}$ & (R) & Can be confirmed by methylene blue test\\\hline
(IV) & NO$_3^-$ & (S) & On reaction with conc. H$_2$SO$_4$, oxidation takes place\\\hline
 & & (T) & All salts are soluble\\\hline
 & & (U) & On reaction with dilute H$_2$SO$_4$ produces colourless gas\\\hline
\end{tabular}}\par
Which of the following is the \textbf{INCORRECT} combination for List-I and List-II.
\begin{qoptions}
\opt{A}{II, S}
\opt{B}{IV, Q}
\opt{C}{I, Q}
\opt{D}{III, U}
\end{qoptions}
\end{bankq}
\begin{bankq}[id=EPS02-P2-CHEM-S3-Q03, source={Enthusiast Pre-Score - Test 02 (11 Nov 2020)}, paper=2, qno=3, type=numerical, subtopic={Anion tests}, answer={4.00}, answer_src=key, difficulty=4, diff_parts={K1 S2 I1 C0 F1}, diff_note={concepts: anions giving coloured gas | steps: reaction with acid; colours; count}, also_in={TOAS - Test 17 (18 Sep 2022) P1 Q3}]
Total number of anions which will give coloured gas/vapour will be :-\par
\bfig{209.6}{EPS02_p2_chem_s3_q03_1.png}
$\mathrm{NO_2^{\ominus}},\mathrm{S_2O_3^{2-}},\mathrm{Cl^-},\mathrm{Br^-},\mathrm{I^-},\mathrm{NO_3^-},\mathrm{F^-}$
\end{bankq}

\section*{TOAS papers}
\begin{bankq}[id=TOAS11-P1-CHEM-S2-Q12, source={TOAS - Test 11 (27 Jun 2022)}, paper=1, qno=12, type=multi, subtopic={Anion tests}, answer={B,D}, answer_src=key, difficulty=2, diff_parts={K1 S1 I0 C0 F1}, diff_note={concepts: tests for anions with dilute sulphuric acid | steps: sulphite and carbonate release gases}]
Which of the following anionic radicals are identified by diluted H$_2$SO$_4$ :--
\begin{qoptions}
\opt{A}{NO$_3^-$}
\opt{B}{SO$_3^{2-}$}
\opt{C}{Cl$^-$}
\opt{D}{CO$_3^{2-}$}
\end{qoptions}
\end{bankq}
\begin{bankq}[id=TOAS12-P1-CHEM-S1-Q09, source={TOAS - Test 12 (10 Jul 2022)}, paper=1, qno=9, type=multi, subtopic={Anion tests}, answer={A,B}, answer_src=key, difficulty=3, diff_parts={K1 S2 I0 C0 F1}, diff_note={concepts: tests for sulphide and nitrite | steps: nitroprusside test; reduction of sulphite by H2S; zinc and KOH on nitrite; brown ring or azo dye test}]
Which of the following statement(s) is/are \textbf{INCORRECT}?
\begin{qoptions}
\opt{A}{H$_2$S when passed in sodium nitroprussiden solution then purple complex is formed}
\opt{B}{H$_2$S reacts with H$_2$SO$_3$ solution and produce sulphur , in this reaction H$_2$S acts as oxidizing agent}
\opt{C}{NaNO$_2$ produce a basic gas on reaction with Zn dust + KOH solution}
\opt{D}{NaNO$_2$ give red coloured product when treated with sulphanilic acid and dil. CH$_3$COOH followed by the addition of 1-nepthylamine.}
\end{qoptions}
\end{bankq}
\begin{bankq}[id=TOAS12-P2-CHEM-S1-Q01, source={TOAS - Test 12 (10 Jul 2022)}, paper=2, qno=1, type=single, subtopic={Anion tests}, answer={B}, answer_src=key, difficulty=2, diff_parts={K1 S2 I0 C0 F0}, diff_note={concepts: reactions of sulphur species | steps: A reduces Fe3+ and dichromate so it is SO2 or H2S; Zn with acid reduces S species; identify A, B, C}]
Consider the following reactions :-\par
A(g) $\xrightarrow{\text{FeCl}_3}$ (B)$\downarrow$ $+$ coloured ion\par
A(g) $\xrightarrow{\text{Cr}_2\text{O}_7^{2-}/\text{H}^{\oplus}}$ (B)$\downarrow$ $+$ coloured ion\par
C(aq.) $\xrightarrow{\text{Zn}+\text{dil. H}_2\text{SO}_4}$ A(g)\par
Identify A, B \& C respectively.
\begin{qoptions}
\opt{A}{SO$_2$, H$_2$S, S$_2$O$_3^{-2}$}
\opt{B}{H$_2$S, S$_8$, SO$_2$}
\opt{C}{SO$_2$, S$_8$, S$^{-2}$}
\opt{D}{SO$_2$, S$_8$, SO$_4^{-2}$}
\end{qoptions}
\end{bankq}
\begin{bankq}[id=TOAS12-P2-CHEM-S3-Q06, source={TOAS - Test 12 (10 Jul 2022)}, paper=2, qno=6, type=numerical, subtopic={Anion tests}, answer={2.00}, answer_src=key, difficulty=3, diff_parts={K1 S2 I0 C0 F1}, diff_note={concepts: redox reactions of anions | steps: anions that decolourise permanganate; anions giving coloured volatile products with HCl; average}]
Consider the following anions?\par
CO$_3^{-2}$, SO$_3^{-2}$, S$^{-2}$, S$_2$O$_3^{-2}$, NO$_2^-$, HCO$_3^-$, HSO$_3^-$, CH$_3$COO$^-$\par
If total number of anion which produce colourless solution with KMnO$_4$/H$^{\oplus}$ by redox reaction $=A$\par
If total number of anion which produce coloured volatile product with dil. HCl $=B$\par
Then find $\dfrac{A+B}{2}=$?
\end{bankq}
\begin{bankq}[id=TOAS13-P2-CHEM-S2-Q09, source={TOAS - Test 13 (24 Jul 2022)}, paper=2, qno=9, type=numerical, subtopic={Anion tests}, answer={4}, answer_src=solved, working={Colourless gas with dil. $\mathrm{H_2SO_4}$: $\mathrm{CO_3^{2-}}$ and $\mathrm{HCO_3^-}$ give $\mathrm{CO_2}$, $\mathrm{S^{2-}}$ gives $\mathrm{H_2S}$, $\mathrm{SO_3^{2-}}$ gives $\mathrm{SO_2}$; $\mathrm{Cl^-}$, $\mathrm{NO_3^-}$ and $\mathrm{BO_3^{3-}}$ give no gas with the dilute acid: 4}, difficulty=2, diff_parts={K1 S1 I0 C0 F1}, diff_note={concepts: anion tests with dilute acid | steps: carbonate, bicarbonate, sulfide and sulfite give colourless gases; the rest do not}]
Find the number of acidic radical(s) which can form colourless gas when treated with dilute H$_2$SO$_4$ ?\par
CO$_3^{2-}$, S$^{2-}$, Cl$^-$, SO$_3^{2-}$, HCO$_3^-$, NO$_3^-$, BO$_3^{3-}$
\end{bankq}
\begin{bankq}[id=TOAS15-P1-CHEM-S1-Q04, source={TOAS - Test 15 (22 Aug 2022)}, paper=1, qno=4, type=single, subtopic={Cation tests}, answer={D}, answer_src=key, difficulty=2, diff_parts={K1 S2 I0 C0 F0}, diff_note={concepts: qualitative tests | steps: nitroprusside needs sulfide ions; chromyl chloride needs ionic chlorides; flame tests; Ca colour through blue glass}]
Which of the following is INCORRECT?
\begin{qoptions}
\opt{A}{H$_2$S does not give test with sodium nitroprusside}
\opt{B}{HgCl$_2$ does not give chromyl chloride}
\opt{C}{Mg$^{2+}$ can not be detected by flame test}
\opt{D}{Through blue glass brick red flame is observe for Ca$^{+2}$ in flame test}
\end{qoptions}
\end{bankq}
\begin{bankq}[id=TOAS15-P2-CHEM-S1-Q05, source={TOAS - Test 15 (22 Aug 2022)}, paper=2, qno=5, type=multi, subtopic={Anion tests}, answer={B}, answer_src=solved, working={A is $\mathrm{Na_2SO_3}$: $\mathrm{SO_2}$ (B) is oxidised by dichromate to sulfate (C), giving $\mathrm{BaSO_4}$ (D); only B is correct among the statements}, difficulty=5, diff_parts={K2 S2 I1 C0 F1}, diff_note={concepts: sulfite tests, dichromate oxidation | steps: SO2 from the dilute acid; oxidised by dichromate to sulfate; precipitates of A with calcium and of D with barium}]
\bfig{187.1}{TOAS15_p2_chem_s1_q05_1.png}
Identify the correct statement :-
\begin{qoptions}
\opt{A}{Salt A having $\text{S}_2\text{O}_3^{2-}$}
\opt{B}{White ppt. (D) is BaSO$_4$}
\opt{C}{B is H$_2$S}
\opt{D}{C form yellow ppt. with HgCl$_2$}
\end{qoptions}
\end{bankq}
\begin{bankq}[id=TOAS16-P1-CHEM-S1-Q01, source={TOAS - Test 16 (04 Sep 2022)}, paper=1, qno=1, type=single, subtopic={Anion tests}, answer={C}, answer_src=key, difficulty=2, diff_parts={K1 S2 I0 C0 F0}, diff_note={concepts: salt analysis statements | steps: chromyl chloride is non-redox; brown ring test interferences; chromate dichromate interconversion needs pH not temperature; borate test}]
Which of the following statement is incorrect?
\begin{qoptions}
\opt{A}{In chromyl chloride test, formation of chromyl chloride is a non redox change}
\opt{B}{The brown ring test of NO$_2^-$ can be done in presence of NO$_3^-$ but brown ring test of NO$_3^-$ cannot be done in presence of NO$_2^-$}
\opt{C}{CrO$_4^{2-}$(aq.) and Cr$_2$O$_7^{2-}$(aq.) can be inter changed by modifying the temperature}
\opt{D}{Borate can be detected by using CaF$_2$ along with conc. H$_2$SO$_4$ without touching the flame}
\end{qoptions}
\end{bankq}
\begin{bankq}[id=TOAS16-P1-CHEM-S1-Q02, source={TOAS - Test 16 (04 Sep 2022)}, paper=1, qno=2, type=single, subtopic={Cation tests}, answer={B}, answer_src=key, difficulty=2, diff_parts={K1 S2 I0 C0 F0}, diff_note={concepts: tests for ammonia | steps: Nessler brown; copper and nickel ammines; neutralising NaOH does not detect ammonia}]
Which of the following observations do not confirm presence of NH$_3$?
\begin{qoptions}
\opt{A}{Nessler's reagent gives brown ppt}
\opt{B}{Neutralising NaOH solution}
\opt{C}{Turning green solution of NiCl$_2$ to deep blue}
\opt{D}{Turning CuSO$_4$ to deep blue}
\end{qoptions}
\end{bankq}
\begin{bankq}[id=TOAS16-P1-CHEM-S1-Q03, source={TOAS - Test 16 (04 Sep 2022)}, paper=1, qno=3, type=single, subtopic={Anion tests}, answer={A}, answer_src=key, difficulty=2, diff_parts={K1 S2 I0 C0 F0}, diff_note={concepts: redox or non-redox nature of reactions | steps: chloride with conc. H2SO4 gives HCl non-redox; bromide is oxidised; thiosulfate reactions are redox}]
Which of the following reaction is non-redox?
\begin{qoptions}
\opt{A}{Cl$^-$ + conc. H$_2$SO$_4\to$}
\opt{B}{Br$^-$ + conc. H$_2$SO$_4\to$}
\opt{C}{Fe$^{3+}$ + S$_2$O$_3^{2-}$ $\xrightarrow[\text{standing}]{\text{on long}}$}
\opt{D}{Cu$^{2+}$ + S$_2$O$_3^{2-}\to$}
\end{qoptions}
\end{bankq}
\begin{bankpassage}[id=TOAS16-P1-CHEM-S3-G1-2, source={TOAS - Test 16 (04 Sep 2022)}, questions={1-2}]
In qualitative analysis, the observations of each test must be carefully observed along with pH conditions and temperature conditions of the solution.
\end{bankpassage}
\begin{bankq}[id=TOAS16-P1-CHEM-S3-Q01, source={TOAS - Test 16 (04 Sep 2022)}, paper=1, qno=1, type=numerical, subtopic={Anion tests}, answer={3}, answer_src=solved, working={Wrong statements: (i) FeSO4 must be freshly prepared, (iii) the test uses freshly made FeSO4 and conc. H2SO4 added slowly down the wall, (vi) a white ppt on adding the acid needs a remedy; (ii), (iv), (v) are right: 3}, difficulty=3, diff_parts={K1 S2 I0 C0 F1}, diff_note={concepts: brown ring test | steps: freshly prepared reagent; conc acid added slowly; interference handling}, group=TOAS16-P1-CHEM-S3-G1-2]
How many of the following steps/observations is/are wrong regarding brown ring test of NO$_3^-$ ?\par
(i) Ferrous sulphate can be allowed to stand before mixing it with the solution\par
(ii) A dark brown ring of nitroso ferrous solution is formed at the junction of the solution\par
(iii) In 1 ml of aqueous solution of salt, 2 mL of dil H$_2$SO$_4$ needs to be added\par
(iv) In case of presence of NO$_2^-$ along with NO$_3^-$, use of sulphamic acid can help with the test.\par
(v) The complex formed is coloured due to charge transfer spectrum\par
(vi) In case of formation of white ppt. while mixing conc. H$_2$SO$_4$, the test can be carried out without any measures.
\end{bankq}
\begin{bankq}[id=TOAS16-P1-CHEM-S3-Q02, source={TOAS - Test 16 (04 Sep 2022)}, paper=1, qno=2, type=numerical, subtopic={Anion tests}, answer={1}, answer_src=solved, working={Nitroprusside $[\mathrm{Fe(CN)_5NO}]^{2-}$ becomes $[\mathrm{Fe(CN)_5NOS}]^{4-}$: charge difference 2; the ambidentate NOS ligand (N- or S-bonded) gives 2 linkage isomers, so the ratio is 1}, difficulty=5, diff_parts={K2 S2 I1 C0 F1}, diff_note={concepts: nitroprusside test, linkage isomerism | steps: charges of reactant and product complexes; difference; linkage isomers of the ambidentate ligand}, group=TOAS16-P1-CHEM-S3-G1-2]
For the complex formed in sodium nitroprusside test of sulphide ion,\par
\par\noindent\hspace*{0.5em}$=\dfrac{\parbox{0.8\linewidth}{\centering\footnotesize Difference between magnitude of charge on product complex ion and reactant complex ion}}{\parbox{0.8\linewidth}{\centering\footnotesize Number of linkage isomers if ambidentate behaviour of CN$^-$ is not considered}}$
\end{bankq}
\begin{bankq}[id=TOAS16-P1-CHEM-S4-Q07, source={TOAS - Test 16 (04 Sep 2022)}, paper=1, qno=7, type=numerical, subtopic={Cation tests}, answer={6}, answer_src=key, difficulty=3, diff_parts={K2 S1 I0 C0 F1}, diff_note={concepts: coloured species, tests | steps: check each of the ten processes for a blue colour; count}]
How many of the following processes involves blue colouration?\par
(1) Cr$_2$O$_7^{2-}$ + H$_2$O$_2$ $\xrightarrow{\text{amyl alcohol}}$\par
(2) Fe$^{2+}$ + [Fe(CN)$_6$]$^{4-}\to$\par
(3) NH$_3$(excess) + NiCl$_2$(aq.)$\to$\par
(4) CuCl$_2$ dil solution in water\par
(5) Green solution of K$_2$MnO$_4$ is presence of acid\par
(6) Cu + ZnSO$_4$ (aq.)$\to$\par
(7) [Co(H$_2$O)$_6$]$^{2+}$ + Cl$^-\to$\par
(8) Fe$^{3+}$ + [Fe(CN)$_6$]$^{4-}\to$\par
(9) NO$_2$ + NO$_2$ $\xrightarrow{\text{excess cooling}}$\par
(10) Al$_2$(SO$_4$)$_3$ + Na$_2$CO$_3$ + Co(NO$_3$)$_2$ $\xrightarrow[\text{fusion in absence of O}_2]{\Delta}$
\end{bankq}
\begin{bankq}[id=TOAS16-P2-CHEM-S1-Q01, source={TOAS - Test 16 (04 Sep 2022)}, paper=2, qno=1, type=multi, subtopic={Cation tests}, answer={A,C}, answer_src=key, difficulty=3, diff_parts={K1 S2 I0 C0 F1}, diff_note={concepts: qualitative analysis statements | steps: nickel bead colour; ammonia dissolves AgCl only; Prussian blue not green; hypo heating}]
Which of the following statement(s) is/are incorrect?
\begin{qoptions}
\opt{A}{Nickel salts gives green borax bead test in oxidising flame.}
\opt{B}{From mixture of AgCl and AgI, NH$_4$OH solution dissolves only AgCl}
\opt{C}{Fe$^{3+}$ ions given a deep green precipitate, on adding K$_4$[Fe(CN)$_6$] solution}
\opt{D}{On strong heating Na$_2$S$_2$O$_3$.5H$_2$O, a polysulphide is formed}
\end{qoptions}
\end{bankq}
\begin{bankpassage}[id=TOAS16-P2-CHEM-S2-G7-8, source={TOAS - Test 16 (04 Sep 2022)}, questions={7-8}]
Two cations on reacting with group-2 reagent gives two different colour ppt. `X' and `Y'.\par
\end{bankpassage}
\begin{bankq}[id=TOAS16-P2-CHEM-S2-Q07, source={TOAS - Test 16 (04 Sep 2022)}, paper=2, qno=7, type=single, subtopic={Cation tests}, answer={A}, answer_src=key, difficulty=2, diff_parts={K1 S2 I0 C0 F0}, diff_note={concepts: group II cations | steps: HCl and water give BiOCl; insoluble in tartaric acid; sulfide colour}, group=TOAS16-P2-CHEM-S2-G7-8]
X + conc. HCl $\xrightarrow{\Delta}$ W + gas 1\par
W + H$_2$O $\rightleftharpoons$ R (white ppt.)\par
R is not soluble in tartaric acid, colour of precipitate `X' is
\begin{qoptions}
\opt{A}{Black}
\opt{B}{Yellow}
\opt{C}{Orange}
\opt{D}{Green}
\end{qoptions}
\end{bankq}
\begin{bankq}[id=TOAS16-P2-CHEM-S2-Q08, source={TOAS - Test 16 (04 Sep 2022)}, paper=2, qno=8, type=single, subtopic={Cation tests}, answer={C}, answer_src=key, difficulty=1, diff_parts={K1 S1 I0 C0 F0}, diff_note={concepts: group II cations | steps: SbOCl dissolves in tartaric acid; sulfide colour}, group=TOAS16-P2-CHEM-S2-G7-8]
Y + conc. HCl $\xrightarrow{\Delta}$ Z + Gas 1\par
Z + H$_2$O $\rightleftharpoons$ Q (white ppt.)\par
Q is soluble in tartatic acid, colour of precipitate `Y' is :-
\begin{qoptions}
\opt{A}{Black}
\opt{B}{Yellow}
\opt{C}{Orange}
\opt{D}{Green}
\end{qoptions}
\end{bankq}
\begin{bankq}[id=TOAS16-P2-CHEM-S4-Q07, source={TOAS - Test 16 (04 Sep 2022)}, paper=2, qno=7, type=numerical, subtopic={Cation tests}, answer={3}, answer_src=key, difficulty=2, diff_parts={K1 S1 I0 C0 F1}, diff_note={concepts: coloured precipitates | steps: find the yellow precipitates among the eight reactions}]
How many of the following will have yellow ppt. as the product.\par
(i) Pb$^{2+}$ (aq.) + 2NaOH $\to$\par
(ii) PbCrO$_4$ + dil. HNO$_3\to$\par
(iii) Ag$_2$CrO$_4$ + excess NH$_3\to$\par
(iv) Hg$^{2+}$ (aq.) + NaOH $\to$\par
(v) K$_2$[Cd(CN)$_4$] + H$_2$S $\to$\par
(vi) Ni$^{2+}$ (aq.) + CN$^-$ (excess) $\to$\par
(vii) H$_2$SO$_4$ (aq.) + (NH$_4$)$_2$CrO$_4\to$\par
(viii) Pb$^{+2}$ (aq.) + (NH$_4$)$_2$CrO$_4\to$
\end{bankq}
\begin{bankq}[id=TOAS17-P1-CHEM-S1-Q06, source={TOAS - Test 17 (18 Sep 2022)}, paper=1, qno=6, type=multi, subtopic={Anion tests}, answer={B,D}, answer_src=key, difficulty=5, diff_parts={K2 S2 I1 C0 F1}, diff_note={concepts: brown ring test, nitroprusside complex | steps: formula and oxidation state of each complex; magnetism; hybridisation; alternative preparation}]
In the following reactions. The correct statements is/are :\par
I. FeSO$_{4\,(\text{aq.})}$ + NO + H$_2$O $\to$ X\par
II. Sodium nitroprusside$_{(\text{aq.})}$ + Na$_2$S $\to$ Y
\begin{qoptions}
\opt{A}{Products X and Y both are paramagnetic}
\opt{B}{In the reaction I change of oxidation state of central atom occurs while in reaction II there is no change in oxidation state of metal}
\opt{C}{Product (X) can also be obtained when NO$_2^-$ or NO$_3^-$ react with dil. H$_2$SO$_4$ in presence of FeSO$_4$}
\opt{D}{Hybridisation of central atom of X is sp$^3$d$^2$ and Y is d$^2$sp$^3$}
\end{qoptions}
\end{bankq}
\begin{bankq}[id=TOAS18-P1-CHEM-S2-Q08, source={TOAS - Test 18 (30 Oct 2022)}, paper=1, qno=8, type=single, subtopic={Cation tests}, answer={A}, answer_src=solved, working={$\mathrm{PbI_2}$ and AgI are yellow, $\mathrm{Hg_2I_2}$ is green, $\mathrm{HgI_2}$ is scarlet red}, difficulty=1, diff_parts={K1 S1 I0 C0 F0}, diff_note={concepts: colours of iodides | steps: colour of each iodide; match the option}]
\begin{center}\footnotesize\setlength{\tabcolsep}{2pt}
\begin{tabular}{|L{0.6cm}L{1.4cm}|L{0.6cm}L{3.3cm}|}\hline
&List-I&&List-II\\\hline
(P)&PbI$_2$&(1)&Water insoluble compound\\\hline
(Q)&AgI&(2)&Yellow colour\\\hline
(R)&Hg$_2$I$_2$&(3)&Green colour\\\hline
(S)&HgI$_2$&(4)&Scarlet red colour\\\hline
&&(5)&White colour\\\hline
\end{tabular}\end{center}
\begin{qoptions}
\opt{A}{P $\to$ 2;Q $\to$ 2;R $\to$ 3;S $\to$ 4}
\opt{B}{P $\to$ 4;Q $\to$ 3;R $\to$ 2;S $\to$ 2}
\opt{C}{P $\to$ 5;Q $\to$ 3;R $\to$ 2;S $\to$ 3}
\opt{D}{P $\to$ 5;Q $\to$ 3;R $\to$ 2;S $\to$ 1}
\end{qoptions}
\end{bankq}
\begin{bankq}[id=TOAS18-P2-CHEM-S2-Q05, source={TOAS - Test 18 (30 Oct 2022)}, paper=2, qno=5, type=multi, subtopic={Anion tests}, answer={A,B,C,D}, answer_src=key, difficulty=2, diff_parts={K1 S1 I0 C0 F1}, diff_note={concepts: chromyl chloride test | steps: only ionic chlorides give the test}]
Which of the following will not give chromyl chloride test.
\begin{qoptions}
\opt{A}{CCl$_4$}
\opt{B}{Hg$_2$Cl$_2$}
\opt{C}{C$_2$H$_5$Cl}
\opt{D}{AgCl}
\end{qoptions}
\end{bankq}
\begin{bankq}[id=TOAS18-P2-CHEM-S3-Q02, source={TOAS - Test 18 (30 Oct 2022)}, paper=2, qno=2, type=single, subtopic={Cation tests}, answer={4}, answer_src=key, difficulty=1, diff_parts={K1 S1 I0 C0 F0}, diff_note={concepts: reaction with hexacyanoferrate(II) | steps: which cations give white precipitates}]
Total number of cation(s) which give white ppt. with K$_4$[Fe(CN)$_6$]
\par\noindent Cu$^{2+}$, Ag$^+$, Ca$^{2+}$, Fe$^{3+}$, Fe$^{2+}$, Zn$^{2+}$, Sr$^{2+}$
\end{bankq}

\section*{TAAS papers}
\begin{bankq}[id=TAAS02-P1-CHEM-S1-Q01, source={TAAS - Test 02 (30 Jul 2023)}, paper=1, qno=1, type=single, subtopic={Anion tests}, answer={C}, answer_src=key, difficulty=2, diff_parts={K1 S2 I0 C0 F0}, diff_note={concepts: identifying a sulfite | steps: gas with a suffocating odour; precipitate with lead acetate; oxidation of the sulfite}]
(I). Salt A $+$ dil. H$_2$SO$_4\rightarrow$ ``B'' gas [B has suffocating order]

(II). Salt A $+$ Pb(OAc)$_2\rightarrow$ C($\downarrow$) $\xrightarrow{\Delta/\mathrm{O_2}}$ D($\downarrow$) [C \& D both are white ppt.]

Have A and D are :-
\begin{qoptions}
\opt{A}{HSO$_3^-$ , PbSO$_4$}
\opt{B}{SO$_3^{2-}$ , PbSO$_3$}
\opt{C}{SO$_3^{2-}$ , PbSO$_4$}
\opt{D}{CO$_3^{2-}$ , PbSO$_3$}
\end{qoptions}
\end{bankq}
\begin{bankq}[id=TAAS02-P1-CHEM-S3-Q08, source={TAAS - Test 02 (30 Jul 2023)}, paper=1, qno=8, type=single, subtopic={Cation tests}, answer={D}, answer_src=key, difficulty=1, diff_parts={K1 S1 I0 C0 F0}, diff_note={concepts: borax bead test colours | steps: colour of each metal in the two flames; match}]
Match the List-I with List-II
\begin{center}\footnotesize\setlength{\tabcolsep}{2pt}
\begin{tabular}{|L{0.6cm}L{1.5cm}L{1.5cm}|L{0.6cm}L{1.0cm}|}\hline
\multicolumn{3}{|L{3.6cm}|}{\centering List-I\\(Colour of borax bead test (cold)}&\multicolumn{2}{L{1.6cm}|}{\centering List-II\\(Metal present)}\\\hline
&Oxidising flame&Reducing flame&&\\\hline
(P)&Blue&Blue&(1)&Cu\\\hline
(Q)&Violet&Colorless&(2)&Fe\\\hline
(R)&Yellow&Green&(3)&Mn\\\hline
(S)&Blue&Red Opaqe&(4)&Co\\\hline
&&&(5)&Cr\\\hline
\end{tabular}\end{center}
The correct option is :
\begin{qoptions}
\opt{A}{P $\to$ 3;Q $\to$ 4;R $\to$ 1;S $\to$ 3}
\opt{B}{P $\to$ 4;Q $\to$ 3;R $\to$ 1;S $\to$ 2}
\opt{C}{P $\to$ 4;Q $\to$ 2;R $\to$ 3;S $\to$ 1}
\opt{D}{P $\to$ 4;Q $\to$ 3;R $\to$ 2;S $\to$ 1}
\end{qoptions}
\end{bankq}
\begin{bankpassage}[id=TAAS02-P2-CHEM-S3-G1-2, source={TAAS - Test 02 (30 Jul 2023)}, questions={1-2}]
Qualitative analysis of inorganic salts the identification of cations and anions present in the salt or a mixture of salts. In the formation of a salt, the part contributed by the acid is known as anion and the part contributed by the base is known as cation.
\bfig{282.6}{TAAS02_p2_chem_s3_para1_2_1.png}
\end{bankpassage}
\begin{bankq}[id=TAAS02-P2-CHEM-S3-Q01, source={TAAS - Test 02 (30 Jul 2023)}, paper=2, qno=1, type=numerical, subtopic={Cation tests}, answer={15}, answer_src=key, difficulty=6, diff_parts={K1 S3 I1 C1 F1}, diff_note={concepts: identification of a mixture from analysis | steps: identify the salts; geometry and hybridisation; isomers; unpaired electrons; evaluate the expression}, group=TAAS02-P2-CHEM-S3-G1-2]
Find value of $(x+2y+3z)$.

where

$x=$ Number of directional orbital used by central atom in hybridisation of anionic part of (G).

$y=$ Total number of geometrical isomers of compound (G).

$z=$ Total number of stereoisomers and unpaired electrons in compound (G).
\end{bankq}
\begin{bankq}[id=TAAS02-P2-CHEM-S3-Q02, source={TAAS - Test 02 (30 Jul 2023)}, paper=2, qno=2, type=numerical, subtopic={Cation tests}, answer={9}, answer_src=key, difficulty=4, diff_parts={K1 S2 I1 C0 F1}, diff_note={concepts: oxidation states in identified compounds | steps: identify H and I; oxidation state of the central atom; zero-valent atoms}, group=TAAS02-P2-CHEM-S3-G1-2]
Find sum of oxidation state of central atom of anion in compound (H) and number of atoms in anion of compound (I) which are in zero oxidation state.
\end{bankq}
\begin{bankq}[id=TAAS02-P2-CHEM-S4-Q05, source={TAAS - Test 02 (30 Jul 2023)}, paper=2, qno=5, type=numerical, subtopic={Anion tests}, answer={1}, answer_src=key, difficulty=2, diff_parts={K1 S1 I0 C0 F1}, diff_note={concepts: inorganic qualitative tests | steps: check each of four statements}]
How many statements is/are correct?

(i) Both PO$_4^{-3}$ \& AsO$_4^{-3}$ give canary yellow ppt. with Ammonium Molybdate under hot \& normal temperature condition respectively.

(ii) Dilute salt of chloride also give chromyl chloride test and it is non-redox reaction

(iii) NO$_3^-$ gives brown ring complex with H$_2$SO$_4$ irrespective of concentration of H$_2$SO$_4$.

(iv) Dichromate ion is more stable in acidic medium and chromate ion is more stable in basic medium
\end{bankq}
\begin{bankq}[id=TAAS02-P2-CHEM-S4-Q06, source={TAAS - Test 02 (30 Jul 2023)}, paper=2, qno=6, type=numerical, subtopic={Anion tests}, answer={1}, answer_src=key, difficulty=2, diff_parts={K1 S1 I0 C0 F1}, diff_note={concepts: inorganic qualitative tests | steps: check each of four statements}]
How many following statements is/are correct?

(i) Except alkali metal carbonate, all carbonates are insoluble in water.

(ii) Cl$_2$ water and H$_2$O$_2$ solution both are used to differentiate CO$_2$ and SO$_2$ gas.

(iii) Iodine Azide test is specific test of Sulphide ion.

(iv) Na$_2$SO$_3$ and NaHSO$_3$ can be differentiated by Litmus paper.
\end{bankq}

\section*{TNAS papers}
\begin{bankq}[id=TNAS06-P2-CHEM-S1-Q02, source={TNAS - Test 06 (30 Jul 2023)}, paper=2, qno=2, type=single, subtopic={Anion tests}, answer={A}, answer_src=key, difficulty=4, diff_parts={K2 S2 I1 C0 F0}, diff_note={concepts: tests for carbonate, sulphide and acetate | steps: solubility; reagents that distinguish; specific tests}]
Which of the following statement is correct?
\begin{qoptions}
\opt{A}{Except alkali metal carbonate and ammonium carbonate , all normal carbonates are insoluble in water}
\opt{B}{Lime water and $\mathrm{Cr_2O_7^{-2}}$ solution both are used to differentiate $\mathrm{CO_2}$ and $\mathrm{SO_2}$ gas.}
\opt{C}{Iodine Azide test is specific test of Sulphide ion only.}
\opt{D}{Acetate ion, gives white ppt. with $\mathrm{FeCl_3}$.}
\end{qoptions}
\end{bankq}
\begin{bankq}[id=TNAS07-P1-CHEM-S1-Q14, source={TNAS - Test 07 (20 Aug 2023)}, paper=1, qno=14, type=single, subtopic={Cation tests}, answer={A}, answer_src=key, difficulty=1, diff_parts={K1 S1 I0 C0 F0}, diff_note={concepts: tests of iron and copper ions with ferricyanide | steps: colour of each precipitate; match}]
{\scriptsize\setlength{\tabcolsep}{2pt}\begin{tabular}{|c|p{0.2\linewidth}|c|p{0.34\linewidth}|}
\hline
\multicolumn{2}{|c|}{\textbf{List-I} Aq. metal ion} & \multicolumn{2}{c|}{\textbf{List-II} Observation with given reagent}\\ \hline
(P) & $\mathrm{Fe^{2+}}$ ion & (1) & blue ppt. with $\mathrm{K_3[Fe(CN)_6]}$\\ \hline
(Q) & $\mathrm{Fe^{3+}}$ ion & (2) & orange ppt. with $\mathrm{K_3[Fe(CN)_6]}$\\ \hline
(R) & $\mathrm{Cu^{2+}}$ ion & (3) & green ppt. with $\mathrm{K_3[Fe(CN)_6]}$\\ \hline
(S) & $\mathrm{Ag^{+}}$ ion & (4) & brown colour with $\mathrm{K_3[Fe(CN)_6]}$\\ \hline
\end{tabular}}\par
\begin{qoptions}
\opt{A}{$P\to1,Q\to4;R\to3;S\to2$}
\opt{B}{$P\to2,Q\to3;R\to1;S\to4$}
\opt{C}{$P\to1,Q\to3;R\to2;S\to4$}
\opt{D}{$P\to4,Q\to1;R\to3;S\to2$}
\end{qoptions}
\end{bankq}
\begin{bankq}[id=TNAS07-P1-CHEM-S2-Q02, source={TNAS - Test 07 (20 Aug 2023)}, paper=1, qno=2, type=single, subtopic={Cation tests}, answer={4}, answer_src=solved, working={$\mathrm{Cu^{2+}}$, $\mathrm{Fe^{3+}}$, $\mathrm{Ni^{2+}}$ and $\mathrm{Mn^{2+}}$ change colour between the flames; $\mathrm{Cr^{3+}}$ is green in both}, difficulty=4, diff_parts={K1 S2 I1 C1 F0}, diff_note={concepts: borax bead colours in two flames | steps: colour in the oxidising flame; colour in the reducing flame; compare; count}]
Number of cations from the following showing different colour in borax bead test in oxidising and reducing flame.\par
$\mathrm{Cu^{2+}}$ , $\mathrm{Fe^{3+}}$ , $\mathrm{Cr^{3+}}$ , $\mathrm{Ni^{2+}}$ , $\mathrm{Mn^{2+}}$
\end{bankq}
\begin{bankq}[id=TNAS07-P2-CHEM-S1-Q01, source={TNAS - Test 07 (20 Aug 2023)}, paper=2, qno=1, type=multi, subtopic={Cation tests}, answer={A,B,C}, answer_src=key, difficulty=7, diff_parts={K1 S3 I2 C1 F1}, diff_note={concepts: identification of a salt from observations | steps: hydrated salt; gases evolved; tests of the solution; identify iron and sulphate; check statements}]
A hydrated green salt (A) loses approx 45\% of its weight on heating and rendered anhydrous. (A) on strong heating yields a dirty brown solid (B) with evolution of two gases (C) and (D). Gas (C) turned orange colour of dichromate solution into green. Dirty brown solid (B) dissolved in boiling HCl solution to produce a yellow solution of (E). Solution (E) was separated into two parts and one part was treated with $\mathrm{NH_4SCN}$ solution giving a blood-red solution. To the other part of solution, $\mathrm{H_2S}$ was passed , where a white turbidity was observed initially and the solution became light green. An aqueous solution of salt (A) gave white precipitate when treated with aqueous barium chloride solution.\par
Consider above given paragraph and select correct option :-
\begin{qoptions}
\opt{A}{Heating of anhydrous salt (A) is an intramolecular redox reaction.}
\opt{B}{C = $\mathrm{SO_2}$}
\opt{C}{D = $\mathrm{SO_3}$}
\opt{D}{Reagent (E) gives blue colour solution with $\mathrm{K_3[Fe(CN)_6]}$}
\end{qoptions}
\end{bankq}

\section*{Star Review (Other) papers}
\begin{bankq}[id=STAR03-P1-CHEM-S4-Q08, source={Star Review (Other) - Test 03 (04 Dec 2023)}, paper=1, qno=8, type=numerical, subtopic={Cation tests}, answer={5}, answer_src=key, difficulty=3, diff_parts={K1 S1 I1 C0 F1}, diff_note={concepts: coloured compounds | steps: colour of each; count green ones}]
Find the number of compound(s) which are green colour .\par
$\mathrm{MnO_2}$, $\mathrm{K_2MnO_4}$, $\mathrm{Cr_2O_3}$, $\mathrm{Ni(OH)_2}$, $\mathrm{CuO}$, $\mathrm{CuCO_3{\cdot}Cu(OH)_2}$, $\mathrm{FeSO_4{\cdot}7H_2O}$
\end{bankq}
\begin{bankq}[id=STAR03-P1-CHEM-S4-Q09, source={Star Review (Other) - Test 03 (04 Dec 2023)}, paper=1, qno=9, type=numerical, subtopic={Cation tests}, answer={1}, answer_src=key, difficulty=4, diff_parts={K1 S2 I1 C0 F1}, diff_note={concepts: solubility of hydroxides in alkali and ammonia | steps: behaviour of each ion with excess alkali; with ammonia; count both}]
Find the number of ions which are soluble in excess of NaOH as well as excess of $\mathrm{NH_4OH}$\par
$\mathrm{Ag^+}$, $\mathrm{Pb^{+2}}$, $\mathrm{Al^{+3}}$, $\mathrm{Sn^{+2}}$, $\mathrm{Zn^{+2}}$, $\mathrm{Cd^{+2}}$.
\end{bankq}
\begin{bankq}[id=STAR03-P2-CHEM-S1-Q06, source={Star Review (Other) - Test 03 (04 Dec 2023)}, paper=2, qno=6, type=multi, subtopic={Anion tests}, answer={A,B,D}, answer_src=key, difficulty=4, diff_parts={K1 S2 I1 C0 F1}, diff_note={concepts: chlorine from chloride with manganese dioxide | steps: reaction; redox; manganese product; its precipitation behaviour}]
When crystals of potassium chloride is heated with $\mathrm{MnO_2}$ and conc. $\mathrm{H_2SO_4}$ then
\begin{qoptions}
\opt{A}{$\mathrm{Cl_2}$ is produced}
\opt{B}{Its redox change}
\opt{C}{Resulting mixture after dilution by water gives buff colour ppt. with $\mathrm{H_2S}$}
\opt{D}{Resulting mixture after dilution by water gives ppt. with excess KOH.}
\end{qoptions}
\end{bankq}

\section*{Nurture (General) papers}
\begin{bankq}[id=NUR01-P2-CHEM-S2-Q04, source={Nurture (General) - Test 01 (12 Jun 2021)}, paper=2, qno=4, type=multi, subtopic={Anion tests}, answer={A,B,C,D}, answer_src=solved, working={CO2, SO2, CO+CO2 and HNO3 vapour}, difficulty=4, diff_parts={K1 S2 I1 C0 F1}, diff_note={concepts: gases from salt and acid reactions | steps: carbonate; sulphite; oxalate; nitrate}]
Which of the following reaction(s) will produce at least one gaseous compound ?
\begin{qoptions}
\opt{A}{$\mathrm{Na_2CO_3(aq)+dil.H_2SO_4}\longrightarrow$}
\opt{B}{$\mathrm{Na_2SO_3(aq)+dil.H_2SO_4}\longrightarrow$}
\opt{C}{$\mathrm{Na_2C_2O_4(aq)+conc.H_2SO_4}\longrightarrow$}
\opt{D}{$\mathrm{NaNO_3(aq)+conc.H_2SO_4}\xrightarrow{\Delta}$}
\end{qoptions}
\end{bankq}
\end{document}
